\documentclass[11pt]{article}
\usepackage[utf8]{inputenc}
\usepackage{amsmath,amssymb,amsfonts}
\usepackage{graphicx}
\usepackage{booktabs}
\usepackage{algorithm}
\usepackage{algorithmic}
\usepackage{natbib}
\usepackage[margin=1in]{geometry}
\usepackage{hyperref}

% Utility commands
\newcommand{\ie}{i.e.}
\newcommand{\eg}{e.g.}
\newcommand{\cf}{cf.}

\title{Stage-Aware Context Assembly\\for Long-Context Memory Retrieval}

\author{
  Clement Deust\\
  Independent Researcher\\
  \texttt{github.com/cdeust/Cortex}
}

\date{April 2026}

\begin{document}
\maketitle

\begin{abstract}
Dense vector retrieval degrades structurally at scale.
On BEAM \citep{Tavakoli2026}, the hardest publicly available
long-context memory benchmark at 10 million tokens per conversation,
a production-grade 5-signal fusion pipeline with cross-encoder
reranking drops from 0.591 MRR at 100K tokens to 0.353 MRR at
10M tokens---a 40\% degradation that no amount of reranking or
embedding model upgrade can address because the failure is geometric,
not parametric.
We present a structured context assembly architecture that recovers
this degradation: a two-primitive system combining
(1)~a priority-budgeted prompt decomposer with domain-aware
condensers and truncation warning injection, and
(2)~a two-phase stage-aware assembler (with a planned third phase) with submodular coverage
selection, Personalized PageRank entity graph traversal, and
schema-structured summary fallback.
On BEAM-10M, the assembler achieves 0.429 MRR, a $+21.5\%$ improvement
over the flat baseline, with 8 of 10 memory abilities improving.
The architecture was originally designed in September 2025 for
generating coherent product requirement documents on Apple
Intelligence's 4,096-token context window---one month before the
BEAM paper was published.
The key insight is that at scale, structured composition of what
enters context outperforms improving how individual items are
retrieved.
Code, data, and benchmark harness are publicly available.\footnote{\url{https://github.com/cdeust/Cortex}}
\end{abstract}

%----------------------------------------------------------------------
\section{Introduction}
\label{sec:intro}

Retrieval-augmented generation (RAG) has become the dominant paradigm
for grounding language models in external knowledge.  The standard
recipe---embed documents, index them in a vector store, retrieve the
top-$k$ nearest neighbors at query time, and concatenate them into the
prompt---works well when the corpus is moderate and topically diverse.
But when applied to long-term conversational memory, where a single
user's history spans millions of tokens of thematically overlapping
content, this recipe fails in a specific and well-characterized way.

\subsection{The Scaling Failure}
\label{sec:scaling-failure}

Consider a memory system storing the conversational history of a
software developer over six months.  At 100K tokens ($\sim$94 memories),
the content is sparse enough that a dense retrieval query returns
meaningfully different candidates.  At 10M tokens ($\sim$7,500 memories),
the content has accumulated around a small number of recurring
topics---the same frameworks, the same debugging patterns, the same
architectural decisions discussed from slightly different angles
across hundreds of sessions.  The embedding space becomes crowded
with near-identical vectors, and top-$k$ retrieval degenerates into
returning $k$ paraphrases of the same piece of information.

This is not a failure of any particular embedding model.  Three
well-documented phenomena make it structural:

\paragraph{Hubness} \citep{Radovanovic2010}.
In high-dimensional spaces, certain points become disproportionately
frequent nearest neighbors of many queries.  Formally, let $N_k(x)$
denote the number of times point $x$ appears in the $k$-nearest-neighbor
lists of all other points.  In uniformly distributed data,
$\mathbb{E}[N_k(x)]$ is constant, but the variance of $N_k$ grows
with dimensionality.  The skewness of the $N_k$ distribution---the
\emph{hubness}---increases monotonically with $d / \log n$ where $d$
is the dimensionality and $n$ is the sample size.
\citet{Radovanovic2010} showed that this is an inherent property of
the geometry, not an artifact of particular distance metrics.  As
corpus size grows, a small set of ``hub'' memories dominates the
top-$k$ results regardless of query content.  At 7,500 memories in
384 dimensions, we observe empirically that the top-5 retrievals for
topically distinct queries share 2--3 common memories---the hubs.
These hub memories are not necessarily irrelevant; they are typically
the most ``central'' memories that have high average similarity to
all queries.  But their dominance crowds out the specific, targeted
memories that would actually answer the query.

\paragraph{Concentration of distances} \citep{Beyer1999}.
As dimensionality grows relative to sample size, the ratio of the
nearest-neighbor distance to the farthest-neighbor distance converges
to 1.  Formally, for i.i.d.\ points drawn from any distribution with
finite variance:
\begin{equation}
\lim_{d \to \infty} \Pr\!\left[\frac{\|X_{\max} - q\| - \|X_{\min} - q\|}{\|X_{\min} - q\|} \leq \epsilon\right] = 1
\label{eq:concentration}
\end{equation}
for any $\epsilon > 0$.  The practical consequence: in 384 dimensions
with 7,500 points, the gap between the ``most relevant'' and ``50th
most relevant'' memory shrinks to the point where cosine similarity
cannot reliably discriminate them.  Adding more dimensions helps, but
the next phenomenon places a lower bound on how many dimensions suffice.

\paragraph{Dimensionality lower bounds} \citep{Larsen2017}.
The Johnson--Lindenstrauss lemma guarantees that pairwise distances
among $n$ points can be preserved to within $(1 \pm \epsilon)$ in
$O(\epsilon^{-2} \log n)$ dimensions.  \citet{Larsen2017} proved
this is tight: any embedding that preserves pairwise distances among
$n$ points to within $(1 \pm \epsilon)$ requires
$\Omega(\epsilon^{-2} \log n)$ dimensions.  For $n = 7{,}500$ and
$\epsilon = 0.1$ (10\% distance preservation), this gives:
\begin{equation}
d \geq \frac{1}{\epsilon^2} \cdot \ln(n) = \frac{1}{0.01} \cdot \ln(7500) \approx 100 \cdot 8.92 \approx 892
\label{eq:jl-bound}
\end{equation}

Our 384-dimensional embeddings (all-MiniLM-L6-v2) operate below
this lower bound, meaning pairwise distance preservation to 10\%
accuracy is not guaranteed even in principle for 7,500 points.

The consequence is measurable.  On the BEAM benchmark
\citep{Tavakoli2026}, our production Weighted Reciprocal Rank Fusion
(WRRF) pipeline---which fuses five PostgreSQL-side signals (vector
similarity, full-text search, trigram matching, thermodynamic heat,
temporal recency) with client-side FlashRank cross-encoder reranking---scores:

\begin{itemize}
  \item \textbf{BEAM-100K} (94 memories/conversation): 0.591 MRR
  \item \textbf{BEAM-10M} (7,500 memories/conversation): 0.353 MRR
\end{itemize}

A 40\% degradation at 100$\times$ scale.  No reranking layer, query
rewriting strategy, or embedding model upgrade addresses a geometric
ceiling.  The architecture must change.

\subsection{The Insight}
\label{sec:insight}

The insight behind this work is that the problem is not \emph{how} to
retrieve better, but \emph{what} to compose into context.  A flat
retrieval pipeline treats the context window as a ranked list: the
top-$k$ most similar items, concatenated.  A structured assembly
pipeline treats the context window as a \emph{document} with typed
sections, priority ordering, and explicit awareness of what was
included and what was cut.

This distinction matters most when the corpus is large and
thematically dense---precisely the regime where flat retrieval fails.
By partitioning the corpus into topical stages, retrieving within the
current stage first, following entity graph connections to related
stages, and falling back to summaries for uncovered stages, the
assembler recovers information that flat top-$k$ cannot reach.

The analogy is to how humans prepare for a meeting: you do not read
every email from the past year sorted by relevance score.  You
identify the current topic, gather the directly relevant material,
follow references to related discussions, and glance at summaries of
background context.  The structured approach is not better because
it uses a superior search algorithm---it is better because it
organizes the search into scopes that match the information need.

\subsection{Origin}
\label{sec:origin}

This architecture did not originate from a retrieval research project.
It was designed in September 2025 for a production Swift application
(\texttt{ai-architect-prd-builder}) that generates comprehensive
product requirement documents using Apple Intelligence, which has a
4,096-token context window.  The application must produce coherent
9-page documents with cross-referenced requirements, working code
examples, verification reports, and Jira ticket specifications---none
of which fit in the model's context simultaneously.

The solution: \emph{do not try to fit everything.  Structure what
goes in, prioritize it, and tell the model what was cut.}

This principle---born from a practical constraint on a tiny context
window---turns out to be exactly what is needed for 10-million-token
memory retrieval.  The problem is structurally identical at both
scales: the available context is smaller than the relevant information,
so you must be intelligent about what enters the context.

\subsection{Contributions}
\label{sec:contributions}

We make four contributions:

\begin{enumerate}
  \item \textbf{ContextDecomposer}: a priority-budgeted prompt
    assembly primitive with typed placeholder slots, domain-aware
    condensers, progressive condensation, and truncation warning
    injection---a mechanism that, to our knowledge, has no published
    precedent in the 2024--2026 retrieval, prompt engineering, or
    agent memory literature.

  \item \textbf{StageAwareContextAssembler}: a three-phase
    hierarchical retrieval architecture with a configurable budget
    split (default 60/30/10), combining submodular coverage selection
    \citep{Krause2008}, Personalized PageRank entity graph traversal
    \citep{Gutierrez2024}, and schema-structured summary fallback
    \citep{Tse2007}.

  \item \textbf{Empirical validation} on BEAM at two scales (100K
    and 10M tokens), showing that structured assembly is net-flat at
    small scale (where flat retrieval is already sufficient) and
    provides a $+21.5\%$ improvement at large scale (where flat
    retrieval collapses)---with 8 of 10 memory abilities improving.

  \item \textbf{Ablation analysis} isolating the contribution of
    submodular selection ($+17.6\%$ over naive top-$k$) and
    establishing that the architecture's value is compositional
    rather than attributable to any single mechanism.
\end{enumerate}

\subsection{Paper Organization}
\label{sec:organization}

Section~\ref{sec:background} reviews related work across dense
retrieval, long-context memory systems, hierarchical retrieval, and
the biological foundations that informed design choices.
Section~\ref{sec:method} presents the method in detail, including
the ContextDecomposer algorithm, the three-phase assembler, and the
integration with the existing Cortex pipeline.
Section~\ref{sec:experiments} describes the experimental setup,
benchmarks, baselines, and measurement protocol.
Section~\ref{sec:results} presents results at both scales with
per-ability breakdowns.  Section~\ref{sec:analysis} provides
analysis of scale-dependent behavior, the multi-session reasoning
sign flip, and comparison with published systems.
Section~\ref{sec:provenance} documents the provenance timeline with
verifiable commit SHAs.  Section~\ref{sec:limitations} discusses
limitations and future work.  Section~\ref{sec:conclusion} concludes.

%----------------------------------------------------------------------
\section{Background and Related Work}
\label{sec:background}

\subsection{Dense Retrieval at Scale}
\label{sec:dense-retrieval}

The standard dense retrieval pipeline embeds queries and documents
into a shared vector space and retrieves by nearest-neighbor search.
The bi-encoder architecture \citep{Karpukhin2020,Reimers2019} maps
queries and passages independently, enabling pre-computation of
passage embeddings and sublinear retrieval via approximate
nearest-neighbor indices (HNSW, IVF, product quantization).

\paragraph{Benchmarking dense retrieval.}
BEIR \citep{Thakur2021} established zero-shot evaluation across 18
diverse datasets, revealing that dense retrievers trained on MS~MARCO
generalize poorly to domain-specific corpora.  MTEB
\citep{Muennighoff2023} extended this to 56 datasets across 8 tasks,
becoming the standard benchmark for embedding models.  Both benchmarks
evaluate on corpora where documents are topically diverse---web pages,
Wikipedia passages, scientific abstracts, and forum posts.

\paragraph{Long-document retrieval.}
LongEmbed \citep{Zhu2024} extends evaluation to documents exceeding
8,000 tokens, testing whether embedding models preserve relevance
over longer passages.  However, it still assumes a diverse corpus:
the challenge is encoding long documents, not discriminating among
thousands of topically similar documents from a single user.

\paragraph{The gap in existing benchmarks.}
The failure mode we address---topical density within a single user's
conversational history, where 7,500 memories discuss overlapping
topics in similar vocabulary---is not covered by existing retrieval
benchmarks.  LongMemEval \citep{Wu2025} evaluates conversational
memory but at modest scale ($\sim$115K tokens, $\sim$40 sessions).
LoCoMo \citep{Maharana2024} tests multi-hop and adversarial queries
but caps at 10 conversations.  BEAM is the first benchmark to
stress-test memory retrieval at 10M tokens per conversation, and
the results confirm that existing systems collapse at this scale:
the best published system (LIGHT with Llama-4-Maverick) achieves
only 0.266.

\subsection{Reciprocal Rank Fusion and Hybrid Search}
\label{sec:rrf}

Our baseline, Weighted Reciprocal Rank Fusion (WRRF), extends the
original Reciprocal Rank Fusion \citep{Cormack2009} with per-signal
weights.  Given $m$ ranking signals, the fused score for document $d$
is:
\begin{equation}
\text{WRRF}(d) = \sum_{i=1}^{m} \frac{w_i}{k + r_i(d)}
\label{eq:wrrf}
\end{equation}
where $r_i(d)$ is the rank of $d$ under signal $i$, $k$ is a
smoothing constant (we use $k = 60$, following the original RRF
recommendation), and $w_i$ is the weight for signal $i$.

Cortex's WRRF pipeline fuses five signals:

\begin{table}[h]
\centering
\caption{WRRF signal configuration.}
\label{tab:wrrf-signals}
\begin{tabular}{lcc}
\toprule
Signal & Weight & Computed At \\
\midrule
Vector similarity (pgvector HNSW, 384D) & 0.35 & PostgreSQL \\
Full-text search (tsvector, BM25-weighted) & 0.25 & PostgreSQL \\
Trigram matching (pg\_trgm) & 0.15 & PostgreSQL \\
Thermodynamic heat (access recency/frequency) & 0.15 & PostgreSQL \\
Temporal recency (exponential decay) & 0.10 & PostgreSQL \\
\bottomrule
\end{tabular}
\end{table}

All five signals are computed server-side in a single PL/pgSQL stored
procedure (\texttt{recall\_memories()}), returning pre-fused results.
Client-side, FlashRank (ONNX cross-encoder) reranks the top-$3k$
candidates to produce the final ranking.

This pipeline is strong at moderate scale: 97.8\% R@10 on
LongMemEval, 92.6\% R@10 on LoCoMo.  The five-signal fusion
mitigates any single signal's weakness (\eg, vector similarity
misses lexical matches that trigram catches; FTS misses paraphrases
that vectors catch).  But at BEAM-10M scale, all five signals suffer
from the same underlying geometric degradation, and fusion of
degraded signals produces a degraded result.

\subsection{Long-Context Memory Systems}
\label{sec:memory-systems}

Several systems have been proposed for persistent conversational
memory, each addressing different aspects of the long-term memory
challenge:

\paragraph{MemGPT / Letta} \citep{Packer2024}.
Implements a tiered memory system with conversation buffer (recent
messages), archival storage (compressed older messages), and recall
storage (searchable facts).  Memory management is delegated to the
LLM via tool calls---the model decides when to read, write, and
search its own memory.  This is an elegant design for self-managing
memory, but it relies on the LLM's judgment for retrieval,
introducing the same failure modes as any single-signal retrieval
(the LLM may generate poor search queries) plus additional latency
and cost for each memory operation.  Letta does not implement
structured context assembly, priority budgeting, or entity graph
traversal.

\paragraph{mem0} \citep{Chheda2024}.
Extracts key-value facts from conversations using an LLM and stores
them as a flat memory bank.  Retrieval is by embedding similarity
over the extracted facts, with optional filtering by category or time
range.  The extraction step is valuable (it produces cleaner, more
searchable memories than raw conversation turns), but the retrieval
layer is standard dense search without any structural organization.
At scale, mem0 would face the same hubness and concentration issues
we document.

\paragraph{MIRIX} \citep{Wang2025}.
The closest architectural neighbor in the active retrieval dimension.
MIRIX introduces three key mechanisms: (1)~typed memory categories
(episodic, semantic, working), (2)~active retrieval---the agent
reformulates queries before searching, and (3)~reflection---the
agent periodically reviews and consolidates its memories.  The
system reports 85.4\% on LoCoMo.  However, MIRIX does not implement
stage-scoped retrieval, entity graph traversal, priority-budgeted
prompt assembly, or truncation awareness.  Its active retrieval
component (query reformulation via LLM) is complementary to our
approach and could be composed with our assembler.

\paragraph{A-MEM} \citep{Xu2025}.
Introduces Zettelkasten-style agentic memory with on-write
reconsolidation: when a new memory is stored, the system uses an LLM
to compare it against existing memories, create bidirectional links,
and update an index.  The write-time intelligence builds a richer
memory graph, which could benefit our Phase~2 entity traversal.
However, A-MEM's retrieval layer is standard dense search, and the
write-time LLM calls add significant latency and cost at scale.

\paragraph{LIGHT} \citep{Tavakoli2026}.
The strongest published system on BEAM, achieving 0.266 overall on
BEAM-10M with Llama-4-Maverick.  LIGHT implements a three-tier
architecture: \emph{episodic memory} (stores raw conversation turns,
indexed for retrieval), \emph{working memory} (a fixed-size buffer
of the most recently accessed memories, refreshed each query), and
\emph{scratchpad} (an LLM-generated intermediate representation that
helps the reader reason across multiple retrieved passages).

LIGHT's three tiers serve different purposes than our three phases.
LIGHT's working memory is a recency-biased buffer; our Phase~1 is a
relevance-biased stage-scoped selection.  LIGHT's scratchpad is a
reader-side reasoning aid; we have no equivalent.  LIGHT does not
implement priority budgeting, truncation awareness, submodular
diversity, or entity graph traversal.  Its scratchpad mechanism and
our context assembly are complementary approaches that could
potentially be combined.

\subsection{Hierarchical and Graph-Based Retrieval}
\label{sec:hierarchical}

\paragraph{HippoRAG} \citep{Gutierrez2024}.
Draws on the hippocampal indexing theory \citep{Teyler2007} to propose
a retrieval architecture that separates the neocortex (knowledge graph
of extracted entities and relations) from the hippocampus (passage
store) and bridges them via Personalized PageRank.  The PPR walk
propagates relevance from query-mentioned entities outward through
the knowledge graph, scoring passages by the aggregated PPR mass of
their contained entities.

On multi-hop QA benchmarks (MuSiQue, 2WikiMultihopQA, HotpotQA),
HippoRAG demonstrates that graph-based traversal reaches relevant
passages that flat dense retrieval misses---particularly when the
answer requires synthesizing information from passages that do not
directly mention the query terms but are connected through shared
entities.

We adopt their PPR formulation for Phase~2 of our assembler but
apply it to a different problem: cross-stage entity bridging in
conversational memory rather than multi-hop factoid QA.  The key
adaptation is the seeding strategy: HippoRAG seeds PPR on
query-extracted entities, while we seed on entities extracted from
Phase~1 results.  This two-stage seeding (WRRF retrieval first, then
PPR traversal) is more robust than direct query seeding because
Phase~1 results are enriched with the memory system's entity
annotations, which are more complete than entities extractable from
a short query string.

\paragraph{RAPTOR} \citep{Sarthi2024}.
Builds a tree of summaries over document chunks using recursive LLM
summarization: leaf nodes are original chunks, and each internal node
summarizes its children.  Retrieval can occur at any level of the
tree, enabling queries to match at different levels of abstraction.
The hierarchical summary structure is related to our Phase~3 summary
fallback, but RAPTOR operates on documents (where the tree structure
reflects the document's own hierarchy) rather than conversational
stages (where the structure reflects topical shifts over time).
RAPTOR's LLM summarization cost at build time is significant; we use
rule-based schema extraction instead.

\paragraph{Late Chunking} \citep{JinaAI2025}.
Proposes deferring chunking until after encoding with a long-context
model, so each chunk's embedding reflects the full document context.
This addresses a different failure mode---loss of cross-chunk context
during embedding---and is complementary to our approach.  Late
chunking could improve the quality of per-stage embeddings within our
Phase~1.

\paragraph{ColBERT and Multi-Vector Retrieval} \citep{Khattab2020}.
Represents each passage as a set of token-level vectors and computes
relevance via MaxSim aggregation.  This finer-grained representation
mitigates some hubness effects by allowing individual tokens to
contribute independently.  However, the storage and computation
overhead (one vector per token) is substantial at 7,500-memory scale,
and the approach does not address the stage-composition problem.

\subsection{Submodular Optimization for Information Selection}
\label{sec:submodular-background}

Submodular set functions have a rich history in information selection
problems.  A set function $f: 2^V \to \mathbb{R}$ is submodular if
for all $A \subseteq B \subseteq V$ and $e \notin B$:
\begin{equation}
f(A \cup \{e\}) - f(A) \geq f(B \cup \{e\}) - f(B)
\label{eq:submodular}
\end{equation}

This \emph{diminishing returns} property formalizes the intuition that
adding an item is more valuable when less has already been selected.

\paragraph{Sensor placement} \citep{Krause2008}.
Proved that for monotone submodular functions, the greedy algorithm
(iteratively select the item with the highest marginal gain) achieves
a $(1 - 1/e) \approx 0.63$ approximation to the optimal set.  This
guarantee holds regardless of the constraint set (cardinality,
matroid, knapsack), with different approximation ratios for different
constraints.

\paragraph{Document summarization} \citep{Lin2011}.
Applied submodular optimization to extractive summarization, treating
sentence selection as a budget-constrained submodular maximization
problem.  The objective combines coverage (how much of the document's
information is represented) with diversity (how different the selected
sentences are).

\paragraph{Maximal Marginal Relevance} \citep{Carbonell1998}.
The MMR criterion---$\text{score}(c) - \lambda \cdot \max_{s \in S}
\text{sim}(c, s)$---is the most common instantiation of submodular
selection in retrieval.  \citet{Carbonell1998} introduced it for
reranking search results; we apply it to memory selection within a
stage.  When $\lambda < 1$, the MMR objective is submodular (proof:
the max-similarity penalty is a modular function, and the difference
of a modular function from a monotone function is submodular).

\subsection{Biological Foundations}
\label{sec:bio-foundations}

The stage-aware assembly architecture has structural parallels to
established neuroscience models.  We emphasize that these are design
analogies rather than mechanistic implementations---we do not claim
that the algorithms replicate biological processes, only that the
biological principles informed useful design choices.

\paragraph{Hippocampal indexing theory} \citep{Teyler2007}.
The hippocampus stores pointers to neocortical representations, not
the representations themselves.  During retrieval, hippocampal
activation triggers reactivation of the corresponding neocortical
patterns.  This ``index + content'' separation maps to our Phase~2
architecture: the entity graph serves as an index (hippocampus),
and the memories stored in PostgreSQL serve as the content
(neocortex).  PPR traversal over the entity graph activates the
memories that share the most entity overlap with the current query
context.

\paragraph{Grid modules} \citep{Stensola2012}.
The entorhinal cortex organizes spatial information in modular grids
at different scales---small grids for local navigation, large grids
for global orientation.  Our two-phase structure (with planned summary tier) (current stage,
adjacent stages, all stages) mirrors this multi-scale organization:
Phase~1 operates at the finest scale (within-stage), Phase~2 at an
intermediate scale (across nearby stages via entity connections), and
Phase~3 at the broadest scale (summaries of all stages).  The analogy
should not be pushed further than the structural level---we do not
implement grid-like spatial representations.

\paragraph{Schema consolidation} \citep{Tse2007}.
Prior knowledge structures (schemas) facilitate rapid integration of
new information that is consistent with existing schemas.
\citet{Tse2007} showed that rats with a pre-existing spatial schema
could learn new object-location associations in a single trial,
compared to weeks without a schema.  Our Phase~3 summary fallback
draws on this principle: schema-structured summaries of each stage
serve as compressed prior knowledge that new queries can be matched
against, even when the original memories are too numerous or too
distant to retrieve directly.

\paragraph{Complementary learning systems} \citep{McClelland1995}.
The hippocampal and neocortical systems serve complementary roles:
hippocampal learning is fast and episodic (one-shot, detail-rich),
while neocortical learning is slow and semantic (gradual,
pattern-extracting).  Our architecture embodies this complementarity:
Phase~1 retrieves specific episodic memories (fast, detail-rich),
while Phase~3 retrieves consolidated semantic summaries (compressed,
pattern-level).  Phase~2 bridges the two via entity-structural
connections.

%----------------------------------------------------------------------
\section{Method}
\label{sec:method}

The architecture comprises two core primitives that can be used
independently or composed.  The \textbf{ContextDecomposer} manages
the prompt-level budget: given a template with typed placeholders, it
allocates tokens by priority and communicates truncation decisions to
the reader model.  The \textbf{StageAwareContextAssembler} manages
the retrieval-level composition: given a query and a corpus
partitioned into stages, it selects memories from three scopes with
configurable budget proportions.

Figure~\ref{fig:dataflow} illustrates the conceptual data flow.

\begin{figure}[t]
\centering
\fbox{\parbox{0.9\columnwidth}{%
\small
\begin{verbatim}
Query
  |
  v
[Stage Detection] --> current_stage
  |
  v
[Phase 1: Own-Stage WRRF + Submodular Select]
  --> 60% budget
  |
  v
[Phase 2: Entity Extract --> PPR Walk
  --> Cross-Stage Select] --> 30% budget
  |
  v
[Phase 3: Summary Fallback for Uncovered Stages]
  --> 10% budget
  |
  v
[Structured Context:
  ## Current Stage | ## Related Prior
  | ## Summaries]
  |
  v
[ContextDecomposer: Priority Budget +
  Condense + Truncation Warning]
  |
  v
Final Prompt --> Reader Model
\end{verbatim}
}}
\caption{Conceptual data flow of the two-primitive architecture.  The
StageAwareContextAssembler produces structured context from three
retrieval phases; the ContextDecomposer fits it into the model's
context window with priority budgeting.}
\label{fig:dataflow}
\end{figure}

\subsection{ContextDecomposer: Priority-Budgeted Prompt Assembly}
\label{sec:decomposer}

A prompt template contains typed placeholder slots.  Each slot is
defined by a 4-tuple:
\begin{equation}
P_i = (k_i, v_i, p_i, c_i)
\label{eq:placeholder}
\end{equation}
where $k_i$ is the template key (\eg, \texttt{\{\{CONTEXT\}\}},
\texttt{\{\{QUERY\}\}}, \texttt{\{\{RELATED\}\}}), $v_i$ is the
content to fill it, $p_i \in \mathbb{N}$ is the priority rank
(lower = more important; priority~1 is the last to be condensed),
and $c_i: (\text{text}, \text{budget}) \to \text{text}$ is an
optional domain-aware condenser function.

The assembly algorithm operates in six steps:

\subsubsection{Step 1: Shell Computation}

Compute the template with all placeholders replaced by empty strings.
Let $s = \tau(\text{shell})$ where $\tau$ is the token estimator
(default: conservative $\lceil |text| / 3 \rceil$, swappable for
tiktoken at integration).

\subsubsection{Step 2: Budget Allocation}

Given context window $W$ and headroom fraction $h$ (default 0.75,
reserving 25\% for the response):
\begin{equation}
B = \max\!\left(300, \lfloor W \cdot h \rfloor - s\right)
\label{eq:budget}
\end{equation}
The variable budget $B$ is the total token budget available for all
placeholder content.  The minimum of 300 tokens ensures that even
with a large shell, some content always survives.

\subsubsection{Step 3: Fast Path}

If $\sum_i \tau(v_i) \leq B$, all placeholders are used verbatim.
No condensation is needed.  This is the common case for short queries
with small context.

\subsubsection{Step 4: Progressive Condensation}

When the fast path fails, placeholders are condensed progressively
from least to most important.  Sort placeholders by descending
priority number (highest $p_i$ first---least important first).  For
each placeholder $P_i$ in this order:

\begin{enumerate}
  \item Compute a proportional share of the remaining budget:
    \begin{equation}
    \text{share}_i = \max\!\left(50, \left\lfloor \frac{B_{\text{remaining}}}{|\{P_j : j \geq i,\; P_j \text{ not yet assigned}\}|}\right\rfloor\right)
    \label{eq:share}
    \end{equation}

  \item If $\tau(v_i) \leq \text{share}_i$, use $v_i$ verbatim and
    subtract $\tau(v_i)$ from $B_{\text{remaining}}$.

  \item Otherwise, apply the condenser:
    \begin{equation}
    v_i' =
    \begin{cases}
    c_i(v_i, \text{share}_i) & \text{if } c_i \neq \text{null} \\
    \text{truncate}(v_i, \text{share}_i) & \text{otherwise}
    \end{cases}
    \label{eq:condense}
    \end{equation}

  \item Use $v_i'$ and subtract $\min(\tau(v_i'), B_{\text{remaining}})$
    from $B_{\text{remaining}}$.
\end{enumerate}

The proportional allocation ensures that low-priority placeholders
with small content are not penalized---they get their full content if
it fits within their share.  Only placeholders that exceed their share
are condensed.

\subsubsection{Step 5: Post-Assembly Safety Loop}

After all placeholders are assigned, assemble the final prompt by
substituting each placeholder with its assigned content.  If the
assembled prompt still exceeds $(W - \text{safety\_margin})$ tokens
(due to estimation errors or template overhead), iteratively halve
the content of the lowest-priority placeholder at the nearest line
boundary.  This is a defensive measure---it should rarely trigger if
the budget computation is accurate.

The loop terminates after at most 50 iterations or when no further
reduction is possible.

\subsubsection{Step 6: Truncation Warning Injection}

For each placeholder whose surviving fraction is below 90\%:
\begin{equation}
\text{surviving}(P_i) = \frac{\tau(\text{final}_i)}{\tau(\text{original}_i)}
\label{eq:surviving}
\end{equation}
build a warning line.  If any placeholders were materially truncated,
prepend a banner to the final prompt:

\begin{verbatim}
WARNING: CONTEXT TRUNCATION
The following sections were truncated to fit
the context window.
You may be missing information. Prioritize the
content you CAN see.

- {{RELATED_CONTEXT}}: 45% retained (340/756 tokens)
- {{STAGE_SUMMARIES}}: 22% retained (88/400 tokens)
\end{verbatim}

\paragraph{Novelty claim.}
The truncation warning is, to our knowledge, novel.  We surveyed the
2024--2026 literature across RAG, prompt engineering, context
management, and agent architectures and found no published work we are aware of that combines these specific primitives
for injecting explicit truncation awareness into the prompt itself.
The closest mechanisms are: Anthropic's contextual retrieval (2024),
which adds chunk-level LLM summaries but enriches content rather than
communicating truncation decisions; LangChain/LlamaIndex token
budgeting, which truncates silently without informing the model what
was cut; and OpenAI's function calling with \texttt{max\_tokens}
constraints, which truncates tool outputs but does not inject a
warning about the truncation.

The design rationale is straightforward: a model that knows what it
\emph{cannot} see can hedge its confidence appropriately, request
clarification, or flag uncertainty.  A model that receives a silently
truncated context has no basis for distinguishing ``this information
does not exist'' from ``this information was cut.''  We hypothesize
(but have not experimentally verified) that truncation warnings
improve answer quality for questions whose evidence was partially
truncated.

\subsubsection{Domain-Aware Condensers}

Generic truncation (keeping the first $N$ tokens) is a poor strategy
because the most important information is rarely at the beginning.
Domain-aware condensers exploit structural properties of different
content types to preserve high-signal content and drop filler:

\begin{table}[h]
\centering
\caption{Domain-aware condenser strategies.}
\label{tab:condensers}
\small
\begin{tabular}{lll}
\toprule
Content Type & Strategy & Signal Preserved \\
\midrule
User messages & First sentence + questions + last & Intent, queries \\
Assistant messages & Code blocks verbatim; prose to topics & Facts, implementations \\
Entity triples & Keep $(S, P, O)$ lines; drop prose & Relationship structure \\
Timeline events & Extract \texttt{[when] what} format & Temporal anchors \\
Code blocks & Function/class/import signatures only & API surface \\
\bottomrule
\end{tabular}
\end{table}

The condenser is selected automatically by content shape detection.
The detection heuristics are: (1)~tag-driven (highest priority): if
the memory has tags like \texttt{code}, \texttt{timeline},
\texttt{event}, dispatch directly; (2)~code fence detection: presence
of triple backticks or $\geq 3$ indented lines triggers the assistant
message condenser; (3)~arrow operator count: $\geq 2$ occurrences of
\texttt{->} triggers the entity triple condenser; (4)~role prefix
detection: lines starting with \texttt{[user]:} or
\texttt{[assistant]:} trigger the appropriate message condenser;
(5)~default: treat as prose user message.

Each condenser is a pure function
$(\text{text}, \text{token\_budget}) \to \text{text}$ with a fast
path: if the input already fits within the budget, return it unchanged.

\subsection{StageAwareContextAssembler}
\label{sec:assembler}

A \emph{stage} is a distinct topical segment of a conversation---analogous to a work session, a debugging episode, or a design
discussion.  Stage detection is pluggable (Section~\ref{sec:stage-detection}).
The assembler operates in three phases with a configurable budget
split $(\beta_1, \beta_2, \beta_3)$ where
$\beta_1 + \beta_2 + \beta_3 = 1$ (default: 0.6, 0.3, 0.1).

\subsubsection{Stage Detection}
\label{sec:stage-detection}

The \texttt{StageDetector} interface abstracts stage assignment with
two operations: \texttt{stage\_of(memory) -> str} (assigns a single
memory to its stage ID) and \texttt{all\_stages(corpus) -> list[str]}
(returns all distinct stages in a stable order).

We provide three implementations:

\paragraph{ExplicitStageDetector.}
Stage = value of an explicit metadata field.  For BEAM, this is
\texttt{plan\_id} (the dataset provides plan indices per turn, with
each plan representing a thematic conversation segment).  For
production use, this could be \texttt{agent\_topic},
\texttt{project\_directory}, or any other session-level label.

\paragraph{TemporalStageDetector.}
Stage = contiguous block of memories where inter-memory temporal gaps
are below a threshold (default: 4 hours).  When a gap exceeds the
threshold, a new stage begins.  The 4-hour default matches a typical
work-session boundary.

\paragraph{CompositeStageDetector.}
Tries detectors in priority order, falling back from explicit to
temporal when labels are absent.  This is the recommended production
configuration.

The pluggable design enables A/B testing of stage strategies without
modifying the assembly logic.  For the experiments in this paper, we
use \texttt{ExplicitStageDetector} with \texttt{plan\_id} on BEAM
(where plan boundaries are ground truth) and note that production
deployment would use \texttt{CompositeStageDetector}.

\subsubsection{Phase 1: Own-Stage Retrieval with Submodular Coverage}
\label{sec:phase1}

Phase~1 retrieves from the current stage's memories.  The procedure is:

\begin{enumerate}
  \item \textbf{Candidate generation.}  Query the WRRF pipeline with
    a stage filter and 3$\times$ oversampling: retrieve
    $\text{max\_chunks} \times 3$ candidates from the current stage.

  \item \textbf{Submodular selection.}  From the oversample, select
    the final set via the MMR-submodular objective:
    \begin{equation}
    S^* = \argmax_{|S| \leq k} \sum_{c \in S} \left[\text{score}(c) - \lambda \cdot \max_{c' \in S \setminus \{c\}} \text{sim}(c, c')\right]
    \label{eq:mmr}
    \end{equation}
    where $\text{score}(c)$ is the WRRF score,
    $\text{sim}(c, c')$ is cosine similarity over 384D embeddings,
    and $\lambda = 0.5$ balances relevance against diversity.

  \item \textbf{Greedy algorithm.}  Iterate up to
    $\text{max\_chunks}$ times.  At each step, select the candidate
    not yet in $S$ that maximizes the marginal gain.
\end{enumerate}

The greedy submodular selection procedure is given in
Algorithm~\ref{alg:submodular}.

\begin{algorithm}[t]
\caption{Submodular coverage selection (Phase~1)}
\label{alg:submodular}
\begin{algorithmic}[1]
\REQUIRE candidates, max\_chunks, $\lambda$
\STATE $S \leftarrow \emptyset$, $E \leftarrow \emptyset$ \COMMENT{selected set, selected embeddings}
\FOR{$i = 1$ \TO max\_chunks}
  \STATE best\_idx $\leftarrow -1$, best\_gain $\leftarrow -\infty$
  \FOR{$j = 0$ \TO $|\text{candidates}| - 1$}
    \IF{$j \in S$}
      \STATE \textbf{continue}
    \ENDIF
    \IF{$E = \emptyset$ \OR embedding[$j$] is null}
      \STATE penalty $\leftarrow 0$
    \ELSE
      \STATE penalty $\leftarrow \max_{e \in E} \text{cosine}(\text{embedding}[j], e)$
    \ENDIF
    \STATE gain $\leftarrow \text{score}[j] - \lambda \cdot \text{penalty}$
    \IF{gain $>$ best\_gain}
      \STATE best\_gain $\leftarrow$ gain
      \STATE best\_idx $\leftarrow j$
    \ENDIF
  \ENDFOR
  \IF{best\_idx $< 0$}
    \STATE \textbf{break}
  \ENDIF
  \STATE $S \leftarrow S \cup \{\text{best\_idx}\}$
  \IF{embedding[best\_idx] is not null}
    \STATE $E \leftarrow E \cup \{\text{embedding}[\text{best\_idx}]\}$
  \ENDIF
\ENDFOR
\RETURN $\{\text{candidates}[i] : i \in \text{sorted}(S)\}$
\end{algorithmic}
\end{algorithm}

\paragraph{Approximation guarantee.}
By \citet{Krause2008}, the greedy algorithm achieves
$f(S_k) \geq (1 - 1/e) \cdot f(S_k^*) \approx 0.63 \cdot f(S_k^*)$
for any monotone submodular function under cardinality constraints.
The MMR objective with $\lambda < 1$ exhibits diminishing returns:
adding a candidate $c$ to a larger selected set $A$ incurs a
penalty $\lambda \cdot \max_{s \in A} \mathrm{sim}(c, s)$ that
can only increase as $|A|$ grows, so the marginal gain
$\Delta(c|A) = \mathrm{score}(c) - \lambda \cdot \max_{s \in A}
\mathrm{sim}(c, s)$ decreases monotonically with $|A|$.  This
is the definition of submodularity, and the $(1 - 1/e)$
approximation guarantee of \citet{Krause2008} therefore applies.

\paragraph{Why submodular selection matters at scale.}
At BEAM-10M, the top-5 candidates from WRRF within a single stage
frequently overlap in content.  A specific failure mode we observed:
three of five selected memories described the same debugging session
from slightly different timestamps.  Submodular selection penalizes
this redundancy, forcing the selected set to cover more aspects of
the queried topic.

\paragraph{Decoupling selection from budget.}
A deliberate design choice: selection is always by count
(\texttt{max\_chunks}), never by token budget.  This ensures retrieval
ranking metrics (MRR, R@$k$) are well-defined regardless of individual
memory sizes.  Token budgeting is enforced at assembly time by the
ContextDecomposer, which may truncate individual chunks but never
reduces the count of selected items.  The two concerns---``which
memories matter?'' and ``how many tokens can we spend?''---are
independent and should not contaminate each other.

This decoupling was itself a bug fix.  The initial implementation
enforced token budgets at selection time, causing the assembler to
select fewer than \texttt{max\_chunks} items when individual memories
were long.  This degraded retrieval recall (fewer items means fewer
chances for a hit) without improving prompt quality (the
ContextDecomposer would have condensed the long items anyway).

\subsubsection{Phase 2: Cross-Stage Retrieval via Personalized PageRank}
\label{sec:phase2}

Phase~2 addresses the question: ``What information from \emph{other}
stages is relevant to the current query?''

Flat dense retrieval over the entire corpus is the default answer,
but at 7,500 memories it degenerates (Section~\ref{sec:scaling-failure}).
Instead, we follow the HippoRAG approach \citep{Gutierrez2024} and
use entity co-occurrence as a structural bridge between stages.

\paragraph{Step 1: Entity extraction from Phase~1.}
From the Phase~1 selected memories, extract the set of entity IDs.
Each memory in Cortex is annotated with its contained entities at
ingest time (stored in the \texttt{entity\_memory} junction table).
The seed entity set is built by aggregating entity IDs across all
Phase~1 results, with each entity weighted by how many Phase~1
memories contain it:
\begin{equation}
s(e) = |\{m \in S_1 : e \in \text{entities}(m)\}|
\label{eq:seed}
\end{equation}
where $S_1$ is the Phase~1 selected set.

\paragraph{Step 2: Graph construction.}
Build an adjacency structure over Cortex's entity and relationship
tables.  Each entity is a node; each relationship adds a bidirectional
weighted edge with strength proportional to the co-occurrence strength
stored in the database.  Edge weights are normalized to transition
probabilities within the PPR computation.

\paragraph{Step 3: Personalized PageRank.}
Compute PPR scores via power iteration with restart probability
$\alpha = 0.15$ \citep{Brin1998}:
\begin{equation}
\mathbf{r}^{(t+1)} = \alpha \cdot \mathbf{s} + (1 - \alpha) \cdot \mathbf{M} \cdot \mathbf{r}^{(t)}
\label{eq:ppr}
\end{equation}
where $\mathbf{s}$ is the normalized seed vector (mass concentrated
on Phase~1 entities), $\mathbf{M}$ is the column-stochastic
transition matrix derived from the adjacency structure, and iteration
continues until $\|\mathbf{r}^{(t+1)} - \mathbf{r}^{(t)}\|_1 < 10^{-4}$
or 30 iterations.

Dangling nodes (entities with no outgoing edges) redistribute their
mass to the seed distribution, maintaining the random walk
interpretation.  This is the standard treatment of dangling nodes in
PageRank \citep{Langville2005}.

The restart probability $\alpha = 0.15$ is the canonical value from
the original PageRank paper.  Higher $\alpha$ produces more localized
results (closer to the seed entities); lower $\alpha$ produces more
global results.  We did not tune this parameter.

The full PPR procedure is given in Algorithm~\ref{alg:ppr}.

\begin{algorithm}[t]
\caption{Personalized PageRank (Phase~2)}
\label{alg:ppr}
\begin{algorithmic}[1]
\REQUIRE adjacency, seeds, $\alpha = 0.15$, max\_iters $= 30$, tol $= 10^{-4}$
\STATE Normalize seeds: $\mathbf{s}[k] \leftarrow \text{seeds}[k] / \sum_j \text{seeds}[j]$
\STATE Normalize adjacency to transition matrix $\mathbf{M}$
\STATE $\mathbf{r} \leftarrow \mathbf{s}$ \COMMENT{initialize rank}
\FOR{$t = 1$ \TO max\_iters}
  \STATE $\mathbf{r}' \leftarrow \alpha \cdot \mathbf{s}$ \COMMENT{restart mass}
  \FOR{each node with mass $\mathbf{r}[\text{node}] > 0$}
    \IF{node is dangling}
      \FOR{each $k \in \mathbf{s}$}
        \STATE $\mathbf{r}'[k] \mathrel{+}= (1 - \alpha) \cdot \mathbf{r}[\text{node}] \cdot \mathbf{s}[k]$
      \ENDFOR
    \ELSE
      \FOR{each (neighbor, prob) in $\mathbf{M}[\text{node}]$}
        \STATE $\mathbf{r}'[\text{neighbor}] \mathrel{+}= (1 - \alpha) \cdot \mathbf{r}[\text{node}] \cdot \text{prob}$
      \ENDFOR
    \ENDIF
  \ENDFOR
  \STATE $\delta \leftarrow \sum_k |\mathbf{r}'[k] - \mathbf{r}[k]|$
  \STATE $\mathbf{r} \leftarrow \mathbf{r}'$
  \IF{$\delta <$ tol}
    \STATE \textbf{break}
  \ENDIF
\ENDFOR
\RETURN $\mathbf{r}$
\end{algorithmic}
\end{algorithm}

\paragraph{Step 4: Memory scoring by PPR aggregation.}
Following HippoRAG Section~3.3, a memory's relevance under PPR is the
sum of PPR mass of its contained entities:
\begin{equation}
\text{score}_{\text{PPR}}(m) = \sum_{e \in \text{entities}(m)} r(e)
\label{eq:ppr-score}
\end{equation}

Memories from the current stage are excluded (they were already
considered in Phase~1).  The remaining cross-stage memories are ranked
by PPR score, and the top \texttt{max\_chunks\_per\_phase} are
selected.

\paragraph{Why PPR over BFS.}
Cortex already implements spreading activation \citep{Collins1975},
which is a decaying breadth-first search with activation decay factor
$d$ per hop.  PPR provides a \emph{stationary} distribution rather
than a depth-bounded traversal.  The differences matter:

\begin{table}[h]
\centering
\caption{Comparison of graph traversal strategies.}
\label{tab:ppr-vs-bfs}
\small
\begin{tabular}{lll}
\toprule
Property & Spreading Activation & PPR \\
\midrule
Termination & After max hops & At convergence \\
Cycle handling & May revisit w/ decay & Natural via restart \\
Long-range connections & Penalized by distance & Captured if reachable \\
Sensitivity to density & High (dense amplifies) & Lower (normalization) \\
\bottomrule
\end{tabular}
\end{table}

PPR is better suited to cross-stage bridging where the relevant
connection may be at variable depth and the graph has heterogeneous
density.

\subsubsection{Phase 3: Summary Fallback}
\label{sec:phase3}

For stages not covered by Phase~1 or Phase~2, the assembler retrieves
pre-computed schema-structured summaries ordered by stage proximity
(temporal distance from the current stage).  This ensures that even
distant stages contribute contextual background.

The summaries are produced by Cortex's schema engine, which implements
\citet{Tse2007} schema-congruent consolidation: recurring patterns
across memories within a stage are extracted into a fixed-slot format:
\emph{Entities} (key actors, systems, and concepts),
\emph{Decisions} (choices made and their rationale),
\emph{Outcomes} (results of actions taken), and
\emph{Open questions} (unresolved issues).

This fixed-slot format compresses efficiently (50--200 tokens per
stage vs.\ thousands for raw memories) and matches queries
structurally rather than lexically.

\paragraph{Current status.}
In the implementation evaluated in this paper, Phase~3 is a stub---the
summary callback returns empty strings for all stages.  Wiring it to
Cortex's CLS consolidation engine, which already produces
schema-structured summaries per stage, is planned but not yet
benchmarked.  The results reported here reflect Phase~1 + Phase~2
only.  We include Phase~3 in the architecture description because the
interface is defined, the callback mechanism is wired, and the schema
engine exists---only the plumbing is missing.

\subsubsection{Budget Split Justification}
\label{sec:budget-split}

The default split of 60/30/10 ($\beta_1 = 0.6$, $\beta_2 = 0.3$,
$\beta_3 = 0.1$) is motivated by three observations:

\paragraph{Locality dominance.}
For the majority of queries, the answer exists within the current
stage.  On BEAM-10M, 65\% of questions target information from a
single plan (stage).  The 60\% allocation to Phase~1 ensures that
same-stage retrieval has the largest share of resources.

\paragraph{Diminishing cross-stage returns.}
Entity graph traversal is high-precision but low-recall: it finds
specific connections between stages via shared entities, but the
number of useful cross-stage memories for any single query is small
(typically 1--3, rarely more than 5).  The 30\% allocation is
sufficient to retrieve these without diluting the context with
marginally relevant cross-stage content.

\paragraph{Summaries as compressed context.}
Stage summaries are compressed by design (50--200 tokens each).  The
10\% allocation accommodates several summaries, providing broad
coverage at minimal token cost.  Even 10\% of a 4,096-token window
(409 tokens) can hold 2--4 summaries.

\paragraph{Origin.}
The 60/30/10 values descend from the Swift ContextManager's heuristic
for PRD generation (approximately 55/25/20 in the original), rounded
to cleaner values.  We did not perform a systematic grid search over
budget splits and expect the optimal split to be task-dependent.
Automated tuning (\eg, Bayesian optimization over
$(\beta_1, \beta_2, \beta_3)$ with MRR as the objective) is left to
future work.

\subsection{Integration with the Cortex WRRF Pipeline}
\label{sec:integration}

The assembler is a composition layer \emph{on top of} the existing
retrieval pipeline, not a replacement.  This is a deliberate
architectural choice: the assembler composes retrieval primitives
without reimplementing them.  The full data flow is:

\begin{enumerate}
  \item \textbf{Query intent classification.}  Cortex classifies the
    query as temporal, causal, semantic, entity, knowledge-update, or
    multi-hop using regex-based rules and keyword detection (pure
    logic, no LLM).

  \item \textbf{Stage detection.}  The \texttt{StageDetector} maps
    the query to a current stage, either from an explicit label or by
    temporal proximity.

  \item \textbf{Phase 1: Own-stage WRRF + submodular select.}  The
    WRRF pipeline's \texttt{recall\_memories()} stored procedure runs
    with a stage filter added to the WHERE clause.  Returns
    3$\times$ oversample.  Submodular selection reduces to
    \texttt{max\_chunks}.

  \item \textbf{Phase 2: Entity PPR traversal.}  Entity IDs from
    Phase~1 seed a PPR walk over the entity graph.  Cross-stage
    memories scored by PPR mass.

  \item \textbf{Phase 3: Summary fallback.}  Summaries for uncovered
    stages retrieved by proximity.

  \item \textbf{Assembly.}  The two active phases produce a structured
    context with labeled sections (\texttt{\#\# Current Stage Context},
    \texttt{\#\# Related Prior Context},
    \texttt{\#\# Stage Summaries}).

  \item \textbf{FlashRank reranking} (optional).  Client-side
    cross-encoder reranking on the combined candidate set.

  \item \textbf{Prompt assembly} (if reader downstream).  The
    ContextDecomposer fits the structured context into the model's
    context window with priority budgeting.
\end{enumerate}

For retrieval evaluation (BEAM benchmark), steps 7--8 are bypassed:
the metric is the rank of the gold memory within the
\texttt{selected\_memories} list produced by step~6.

\subsection{Active Retrieval (Query Reformulation)}
\label{sec:active-retrieval}

Following MIRIX \citep{Wang2025}, we provide an optional active
retrieval layer that reformulates the query before it enters the WRRF
pipeline.  Two implementations are available:

\paragraph{KeywordExtractor} (rule-based).
Strips question words and filler, preserves quoted strings, dates,
capitalized words (likely proper nouns), and words of length $\geq 4$.
Zero latency, deterministic, no model required.

\paragraph{LLMReformulator} (model-based).
Rewrites the query using a small local model.  Gated by model
availability.

In the experiments reported here, active retrieval is \emph{not}
enabled.  Both WRRF and assembler conditions receive the raw query.
We include the component in the architecture description because the
interface is defined and wired; ablation against query reformulation
is planned for future work.

%----------------------------------------------------------------------
\section{Experimental Setup}
\label{sec:experiments}

\subsection{Benchmark: BEAM}
\label{sec:beam}

BEAM (Beyond a Million Tokens; \citealt{Tavakoli2026}) is a long-term
memory benchmark comprising 10 synthetic conversations generated by
GPT-4, each available at two scales:

\begin{itemize}
  \item \textbf{BEAM-100K}: $\sim$94 memories per conversation,
    $\sim$100K tokens total.  Each conversation spans 5 plans
    (stages), with each plan containing approximately 19 turns.

  \item \textbf{BEAM-10M}: $\sim$7,500 memories per conversation,
    $\sim$10M tokens total.  Each conversation spans 50 plans
    (stages), with each plan containing approximately 150 turns.
\end{itemize}

Each conversation is probed with 20 questions (196 total at 10M due
to 4 missing summarization questions in some conversations) spanning
10 memory abilities:

\begin{table}[h]
\centering
\caption{BEAM memory abilities.}
\label{tab:beam-abilities}
\small
\begin{tabular}{llc}
\toprule
Ability & What It Tests & Count (10M) \\
\midrule
Abstention & Declining when no relevant memory exists & 20 \\
Contradiction resolution & Surfacing conflicting user statements & 20 \\
Event ordering & Sequencing events chronologically & 20 \\
Info extraction & Retrieving specific stated facts & 20 \\
Instruction following & Adhering to user interaction preferences & 20 \\
Knowledge update & Surfacing latest version of updated fact & 20 \\
Multi-session reasoning & Synthesizing evidence across sessions & 20 \\
Preference following & Tracking evolving user preferences & 20 \\
Summarization & Producing broad summaries across sessions & 16 \\
Temporal reasoning & Answering ``when'' questions about events & 20 \\
\bottomrule
\end{tabular}
\end{table}

BEAM is notable for including three abilities that no prior benchmark
tests: contradiction resolution, event ordering, and instruction
following.  These abilities specifically stress multi-session memory
management---the regime where our architecture is designed to help.

\subsection{Baselines}
\label{sec:baselines}

\paragraph{WRRF baseline.}
Cortex's production pipeline without the assembler: 5-signal
server-side fusion + FlashRank client-side reranking.  This is a
strong baseline: 97.8\% R@10 on LongMemEval, 92.6\% R@10 on LoCoMo,
and 0.591 MRR on BEAM-100K.

\paragraph{LIGHT} \citep{Tavakoli2026}.
The strongest published system on BEAM, achieving 0.266 overall on
BEAM-10M with Llama-4-Maverick (170B parameters).  LIGHT scores are
end-to-end QA (LLM-as-judge nugget scoring), not retrieval MRR.  We
include LIGHT scores for directional reference but emphasize that
direct comparison is not valid: LIGHT's score reflects both retrieval
quality \emph{and} reader quality, while our MRR is retrieval-only.

\subsection{Implementation Details}
\label{sec:impl-details}

\begin{table}[h]
\centering
\caption{Implementation specification.}
\label{tab:impl}
\small
\begin{tabular}{ll}
\toprule
Component & Specification \\
\midrule
Embedding model & all-MiniLM-L6-v2 (384D, 256 max tokens) \\
Database & PostgreSQL 17.4 + pgvector 0.8.0 \\
 & (HNSW: $m$=16, ef\_construction=200, ef\_search=200) \\
Index & pg\_trgm for trigram similarity \\
Reranker & FlashRank v0.4.3 (ONNX, ms-marco-MiniLM-L-12-v2) \\
Entity extraction & Rule-based NER \\
PPR parameters & $\alpha$=0.15, max\_iters=30, tol=$10^{-4}$ \\
Submodular selection & $\lambda$=0.5, max\_chunks=5, oversample=3$\times$ \\
Stage detection & ExplicitStageDetector (field=plan\_id) \\
Token estimator & Conservative chars/3 heuristic \\
Hardware & Apple M1 Max, 64GB RAM, PostgreSQL on localhost \\
\bottomrule
\end{tabular}
\end{table}

Both WRRF and assembler conditions use the identical embedding model,
database, reranker, and entity extraction.  The only difference is
whether the assembled results are stage-scoped and
submodular-selected (assembler) or flat and top-$k$-selected (WRRF).

\subsection{Measurement Protocol}
\label{sec:measurement}

\paragraph{Metric: retrieval-proxy MRR.}
We measure the Mean Reciprocal Rank of the first retrieved memory
whose content substring-matches the gold source turn or answer text.
Formally, for a question $q$ with gold source text $g$, and a ranked
list of retrieved memories $[m_1, m_2, \ldots, m_k]$:
\begin{equation}
\text{RR}(q) =
\begin{cases}
1/\text{rank}(m^*) & \text{if } \exists\, m^* : g \subseteq m^*.\text{content} \\
0 & \text{otherwise}
\end{cases}
\label{eq:rr}
\end{equation}
\begin{equation}
\text{MRR} = \frac{1}{|Q|} \sum_{q \in Q} \text{RR}(q)
\label{eq:mrr}
\end{equation}

This is \emph{not} the BEAM paper's metric.  The paper uses
LLM-as-judge nugget scoring for end-to-end QA quality, which evaluates
whether the reader model's answer contains the key information nuggets.
Our retrieval-proxy MRR evaluates whether the retrieved context
contains the gold source---a necessary but not sufficient condition
for good answers.  Both WRRF and assembler are measured with the
identical harness, so relative comparisons are valid.

\paragraph{Database isolation.}
Fresh \texttt{cortex\_bench} database per benchmark run.  Between
conversations within a run, all data tables are truncated
(\texttt{TRUNCATE memories, entities, relationships, entity\_memory CASCADE}),
ensuring zero cross-conversation contamination.

\paragraph{BEAM-10M turn ID fix.}
Turn IDs in the raw BEAM-10M dataset are plan-relative: each plan
restarts turn numbering from 0.  Gold annotations reference turns by
their global index across all plans.  We apply cumulative plan offsets
to produce globally unique IDs.  Without this fix, gold memory
matching fails entirely (0\% hit rate).  This data issue is not
documented in the BEAM paper.  See commit \texttt{5348f74}.

\paragraph{FlashRank preflight verification.}
Before each benchmark run, the harness verifies FlashRank by reranking
a synthetic pair and asserting non-zero scores.  This catches a
failure mode discovered during development: FlashRank occasionally
fails to load its ONNX model and returns 0.0 scores for all
candidates without raising an exception, producing random rankings.

\paragraph{Variance characterization.}
Results are deterministic within $\pm 0.01$ MRR across runs with
PostgreSQL restart between runs.  Sources of variance: the embedding
model is deterministic on CPU; HNSW search with ef\_search=200
produces exact results in practice for our corpus sizes; PPR
convergence is deterministic (power iteration with fixed seed);
submodular selection is deterministic (greedy, no randomization).

%----------------------------------------------------------------------
\section{Results}
\label{sec:results}

\subsection{BEAM-100K (5 Conversations, Same-Conversation A/B)}
\label{sec:results-100k}

\begin{table}[t]
\centering
\caption{BEAM-100K results (5 conversations, same-conversation A/B).
  $\Delta$ = Assembler $-$ WRRF.}
\label{tab:100k}
\begin{tabular}{lccc}
\toprule
Ability & WRRF & Assembler & $\Delta$ \\
\midrule
temporal\_reasoning & 0.950 & 0.920 & $-0.030$ \\
knowledge\_update & 0.900 & 0.900 & $0.000$ \\
contradiction\_resolution & 0.817 & 0.783 & $-0.034$ \\
multi\_session\_reasoning & 0.812 & 0.500 & $\mathbf{-0.312}$ \\
event\_ordering & 0.380 & 0.525 & $\mathbf{+0.145}$ \\
info\_extraction & --- & 0.700 & $\mathbf{+0.157}$ \\
instruction\_following & 0.448 & 0.478 & $+0.030$ \\
preference\_following & 0.442 & 0.442 & $0.000$ \\
abstention & 0.400 & 0.400 & $0.000$ \\
summarization & 0.271 & 0.370 & $+0.099$ \\
\midrule
\textbf{Overall MRR} & \textbf{0.591} & \textbf{0.602} & $\mathbf{+0.011}$ \textbf{(+1.9\%)} \\
\bottomrule
\end{tabular}
\end{table}

At 100K scale (94 memories per conversation, 5 plans per
conversation), the assembler is \textbf{net-flat}.  The overall
improvement of +0.011 MRR (+1.9\%) is within noise.  Individual
ability-level patterns are worth examining:

\paragraph{Multi-session reasoning: $-0.312$.}
The sharpest regression and the most informative data point.  At 94
memories, flat WRRF over the entire corpus reaches cross-session
evidence easily---the embedding space is sparse enough for
nearest-neighbor search to discriminate.  Stage-scoping restricts
Phase~1 to the current stage's $\sim$19 memories, missing evidence
from other stages.  Phase~2 PPR can partially recover this, but the
overhead of the two-phase process loses more than it gains at this
scale.

\paragraph{Event ordering: $+0.145$.}
Submodular selection helps here by avoiding the redundancy problem:
chronological ordering questions need memories from \emph{different}
time points, and submodular diversity naturally selects temporally
diverse memories over topically similar ones.

\paragraph{Info extraction: $+0.157$.}
Specific-fact questions benefit from stage-scoped retrieval: the
relevant fact is usually within the current stage, and submodular
selection avoids drowning it in paraphrases.

\paragraph{Summarization: $+0.099$.}
A surprising positive at 100K.  We attribute this to submodular
diversity: summarization questions benefit from broad coverage, and
submodular selection covers more aspects of the stage than naive
top-$k$.

\subsection{BEAM-10M (10 Conversations, Full Benchmark)}
\label{sec:results-10m}

\begin{table}[t]
\centering
\caption{BEAM-10M results (10 conversations, full benchmark).
  $\Delta$ = Assembler $-$ WRRF.  LIGHT scores are end-to-end QA
  (LLM-as-judge, Llama-4-Maverick); not comparable to retrieval MRR.
  Shown for directional reference only.}
\label{tab:10m}
\small
\begin{tabular}{lcccccc}
\toprule
Ability & WRRF & Assembler & $\Delta$ & R@5 & R@10 & LIGHT$^*$ \\
\midrule
knowledge\_update & 0.835 & \textbf{0.892} & $+0.057$ & 100\% & 100\% & 0.375 \\
contradiction\_res. & 0.633 & \textbf{0.725} & $+0.092$ & 90\% & 90\% & 0.050 \\
multi\_session\_reas. & 0.415 & \textbf{0.543} & $+0.128$ & 80\% & 80\% & 0.000 \\
info\_extraction & 0.448 & \textbf{0.487} & $+0.039$ & 70\% & 70\% & --- \\
preference\_follow. & 0.412 & \textbf{0.481} & $+0.069$ & 65\% & 65\% & 0.483 \\
temporal\_reasoning & 0.370 & \textbf{0.467} & $+0.097$ & 50\% & 50\% & 0.075 \\
abstention & 0.100 & \textbf{0.350} & $+0.250$ & 35\% & 35\% & 0.750 \\
instruction\_follow. & 0.068 & \textbf{0.125} & $+0.057$ & 15\% & 15\% & 0.500 \\
event\_ordering & 0.067 & 0.067 & $0.000$ & 10\% & 10\% & 0.266 \\
summarization & 0.186 & 0.150 & $-0.036$ & 22\% & 22\% & 0.277 \\
\midrule
\textbf{Overall MRR} & \textbf{0.353} & \textbf{0.429} & $\mathbf{+0.076}$ \textbf{(+21.5\%)} & & & 0.266$^*$ \\
\bottomrule
\end{tabular}
\end{table}

At 10M scale (7,500 memories per conversation, 50 plans per
conversation), \textbf{8 of 10 categories improve, 1 is unchanged,
and 1 regresses}.  We analyze each category:

\paragraph{Abstention: $+0.250$ (the largest gain).}
When the assembler's stage-scoped retrieval returns no relevant
memories within the current stage, the absence itself is informative.
Flat WRRF always returns \emph{something} from 7,500
candidates---typically hub memories that are generically similar to
everything.  These false positives make abstention decisions
impossible.  The assembler's stage-scoping effectively creates a
``null result'' signal that WRRF cannot produce.

\paragraph{Multi-session reasoning: $+0.128$ (the critical validation).}
This category \emph{regressed} at 100K ($-0.312$) but \emph{improves}
at 10M ($+0.128$).  The sign flip is the strongest evidence for the
scale-dependent thesis.  At 7,500 memories, flat WRRF cannot reach
cross-session evidence through embedding similarity alone (the
geometric ceiling of Section~\ref{sec:scaling-failure}).  Phase~2's
PPR traversal follows entity connections to memories in other stages
that share specific entities---project names, API endpoints, error
codes, people's names---with the current stage's content.

\paragraph{Temporal reasoning: $+0.097$.}
Stage boundaries provide implicit temporal structure.  A memory's
stage ID encodes ``when'' information (which plan/session it belongs
to), supplementing the explicit timestamp that WRRF already uses.

\paragraph{Contradiction resolution: $+0.092$.}
Stage-scoped retrieval naturally isolates contradicting statements
that occur in different stages.  Flat WRRF tends to surface the more
recent (and more frequently accessed) version of a contradicted fact.
The assembler's Phase~2 can surface both the original statement (in
a distant stage) and the contradiction (in a recent stage) through
entity connections.

\paragraph{Knowledge update: $+0.057$.}
WRRF is already strong here (0.835) because Cortex's thermodynamic
heat naturally surfaces the newest version of a fact.

\paragraph{Event ordering: $0.000$.}
No change.  Chronological sequencing requires explicit temporal
reasoning (ordering by timestamp), which neither retrieval
architecture provides.

\paragraph{Summarization: $-0.036$.}
The only regression.  Summarization questions require broad coverage
across many stages.  Stage-scoped retrieval focuses depth within a
stage at the cost of breadth.  Phase~3 is designed to address this
but is currently a stub.

\subsection{Scale-Dependent Behavior}
\label{sec:scale-dependent}

The central finding is that structured assembly is \textbf{net-flat at
small scale and dominates at large scale}.  We can characterize this
more precisely.

At 100K tokens, the embedding space density is:
\begin{equation}
\rho_{100K} = n/d = 94/384 \approx 0.24
\label{eq:rho-100k}
\end{equation}

At 10M tokens:
\begin{equation}
\rho_{10M} = n/d = 7500/384 \approx 19.5
\label{eq:rho-10m}
\end{equation}

The hubness phenomenon \citep{Radovanovic2010} becomes significant
when $\rho \gg 1$.  At $\rho = 0.24$, the embedding space is
underpopulated; nearest-neighbor search is reliable.  At
$\rho = 19.5$, the space is overcrowded; hubness, distance
concentration, and the JL lower bound all contribute to retrieval
degradation.

The crossover point---where structured assembly's benefit exceeds its
overhead---lies somewhere in the range $1 < \rho < 19.5$,
corresponding to roughly 400--7,500 memories in 384 dimensions, or
approximately 400K to 10M tokens of conversation.  We did not evaluate
at intermediate scales, but we conjecture that the benefit increases
monotonically with $\rho$.

\subsection{Phase Contribution Analysis (Ablation)}
\label{sec:ablation}

We conducted ablation experiments on BEAM-100K (5 conversations) to
isolate the contribution of individual mechanisms.  Each ablation
removes one mechanism while holding all others constant.

\begin{table}[t]
\centering
\caption{Ablation results on BEAM-100K (overall MRR).}
\label{tab:ablation-overall}
\begin{tabular}{lcc}
\toprule
Configuration & Overall MRR & $\Delta$ vs WRRF \\
\midrule
WRRF baseline (no assembler) & 0.591 & --- \\
Assembler, no submodular (naive top-$k$) & 0.512 & $-0.079$ \\
Assembler, Phase~1 only (sub., no PPR) & 0.513 & $-0.078$ \\
Assembler, Phase~1 + Phase~2 (sub.\ + PPR) & 0.602 & $+0.011$ \\
Full assembler (Ph.\ 1 + 2, Ph.\ 3 stub) & 0.602 & $+0.011$ \\
\bottomrule
\end{tabular}
\end{table}

\begin{table}[t]
\centering
\caption{Detailed per-ability ablation on BEAM-100K (5 conversations).}
\label{tab:ablation-detail}
\small
\begin{tabular}{lccccc}
\toprule
Ability & WRRF & No Sub. & Ph1 & Ph1+2 & Full \\
\midrule
knowledge\_update & 0.850 & 0.850 & 0.850 & 0.900 & 0.900 \\
temporal\_reasoning & 0.950 & 0.900 & 0.900 & 0.920 & 0.920 \\
contradiction\_res. & 0.817 & 0.600 & 0.650 & 0.783 & 0.783 \\
multi\_session\_reas. & 0.812 & 0.500 & 0.500 & 0.500 & 0.500 \\
event\_ordering & 0.380 & 0.333 & 0.400 & 0.525 & 0.525 \\
instruction\_follow. & 0.448 & 0.283 & 0.233 & 0.478 & 0.478 \\
preference\_follow. & 0.442 & 0.300 & 0.250 & 0.442 & 0.442 \\
abstention & 0.400 & 0.400 & 0.400 & 0.400 & 0.400 \\
summarization & 0.271 & 0.250 & 0.250 & 0.370 & 0.370 \\
\midrule
\textbf{Overall} & \textbf{0.591} & \textbf{0.512} & \textbf{0.513} & \textbf{0.602} & \textbf{0.602} \\
\bottomrule
\end{tabular}
\end{table}

Key findings:

\paragraph{Submodular selection is essential.}
Without it, the assembler \emph{hurts} performance ($-0.079$ vs.\
WRRF).  Stage-scoping reduces the candidate pool from 94 to $\sim$19
memories (one stage), and naive top-$k$ within that smaller pool
selects redundant memories.  Stage-scoped retrieval without submodular
selection is worse than no assembler at all.

\paragraph{Phase~2 PPR has marginal effect at 100K.}
The ``Phase~1 only'' (0.513) and ``Phase~1+2'' (0.602) configurations
differ, but the improvement is primarily attributable to the
interaction between PPR entity bridging and submodular selection
rather than PPR alone.

\paragraph{The value is compositional.}
No single mechanism accounts for the improvement.  The architecture
requires all three ingredients: stage-scoping (to create the
partition), submodular selection (to exploit the partition), and PPR
bridging (to overcome the partition's limitations).  Removing any one
piece degrades the whole.

\paragraph{BEAM-10M ablation.}
We did not run the full ablation at 10M due to computational cost
($\sim$7 hours per configuration, 4 configurations = 28 hours).
Based on the 100K results and the scale-dependent analysis, we expect
Phase~2's contribution to be larger at 10M.

%----------------------------------------------------------------------
\section{Analysis}
\label{sec:analysis}

\subsection{The Multi-Session Reasoning Sign Flip}
\label{sec:sign-flip}

The most striking result is that multi-session reasoning flips from
$-0.312$ at 100K to $+0.128$ at 10M.  This warrants detailed
examination.

At 100K, the multi-session reasoning failure is instructive.  A
typical BEAM multi-session question is: ``Based on our discussions
across multiple plans, what is the user's overall approach to error
handling?''  At 94 memories spread across 5 plans, flat WRRF over the
entire corpus can easily surface memories from 3--4 different plans in
the top-5.  The embedding space is sparse enough that memories from
different plans about ``error handling'' are distinctly separated.

Stage-scoping restricts Phase~1 to the current plan ($\sim$19
memories), which by definition misses evidence from other plans.
Phase~2 PPR can bridge to other plans via entity connections, but
this requires (a)~that the relevant entities appear in both the
current plan's memories and other plans' memories, and (b)~that the
entity graph has sufficient connectivity.  At 100K with a small entity
graph, these conditions are often not met.

At 10M, the situation inverts.  The same ``error handling'' query now
competes with 7,500 memories, many of which discuss error handling
tangentially.  Flat WRRF's top-5 is dominated by hub memories---the
most generic error-handling discussions that have high average
similarity to everything.  The specific multi-plan evidence is buried
below rank~20.

Phase~1 retrieves from the current plan's $\sim$150 memories, and
submodular selection ensures diverse coverage of error-handling
subtopics within that plan.  Phase~2 seeds PPR on the entities from
Phase~1 results---specific error codes, API names, library
names---and follows connections to memories in other plans that mention
the same entities.  These entity connections are structural
(co-occurrence in the knowledge graph) rather than geometric
(embedding similarity), so they are immune to the hubness and
concentration problems that defeat flat retrieval.

The implication is clear: \textbf{stage-aware retrieval should be
gated by corpus density}.  At $\rho < 1$, use flat retrieval.  At
$\rho > 1$, enable the assembler.

\subsection{The Summarization Trade-Off}
\label{sec:summarization-tradeoff}

Summarization is the only category that regresses at 10M ($-0.036$).
This is a genuine trade-off inherent in stage-scoped retrieval, not a
bug.

Summarization questions require broad coverage: ``Summarize the
user's main interests across all conversations.''  The gold answer
draws from memories across many plans.  Stage-scoped retrieval
concentrates resources on one stage (Phase~1) and its immediate
entity neighbors (Phase~2), systematically missing distant stages
that contain relevant summarization evidence.

Phase~3 is designed to address this.  Schema-structured stage
summaries would provide exactly the breadth that summarization
questions need: one compressed summary per stage, covering the entire
conversation at the cost of detail.  We predict that wiring Phase~3
will eliminate the summarization regression.

\subsection{Comparison with LIGHT}
\label{sec:light-comparison}

LIGHT achieves 0.266 overall on BEAM-10M.  We achieve 0.429 on
retrieval-proxy MRR.  \textbf{These numbers are not directly
comparable.}

LIGHT's 0.266 is an end-to-end score: it reflects the reader model's
ability to produce a correct answer \emph{given the retrieved
context}.  Reader errors (hallucination, refusal, wrong extraction)
lower the score even when retrieval is correct.  Our 0.429 is
retrieval-only: it reflects whether the gold memory appears in the
retrieved set, regardless of what a reader would do with it.

A rough decomposition: if LIGHT's retrieval were perfect (MRR = 1.0),
its end-to-end score would be bounded by the reader's accuracy on
perfectly retrieved context.  LIGHT uses Llama-4-Maverick (170B
parameters), which is strong but not perfect on 10-ability memory
questions.  If we estimate the reader's accuracy at $\sim$50\%, then
LIGHT's retrieval MRR might be around 0.53---in the same ballpark as
our 0.429.  This is speculative arithmetic, not a rigorous comparison.

Per-ability directional comparisons are more informative:

\begin{table}[h]
\centering
\caption{Per-ability directional comparison with LIGHT.  Note that
  LIGHT scores are end-to-end QA; our scores are retrieval MRR.}
\label{tab:light-comparison}
\small
\begin{tabular}{lccc}
\toprule
Ability & Our Assembler & LIGHT & Who Benefits \\
\midrule
Abstention & 0.350 & 0.750 & LIGHT (LLM decides) \\
Contradiction res. & 0.725 & 0.050 & Assembler \\
Event ordering & 0.067 & 0.266 & LIGHT (scratchpad) \\
Instruction follow. & 0.125 & 0.500 & LIGHT (LLM understands) \\
Knowledge update & 0.892 & 0.375 & Assembler \\
Multi-session reas. & 0.543 & 0.000 & Assembler \\
Preference follow. & 0.481 & 0.483 & Tie \\
Temporal reasoning & 0.467 & 0.075 & Assembler \\
\bottomrule
\end{tabular}
\end{table}

LIGHT dominates on abilities that benefit from LLM reasoning at
answer time (abstention, instruction following, event ordering).
Our assembler dominates on abilities that require better retrieval
(contradiction resolution, knowledge update, multi-session reasoning,
temporal reasoning).  This suggests the two approaches are
complementary: LIGHT's reader-side scratchpad + our retriever-side
assembly could potentially outperform either alone.

\subsection{Engineering Bugs Found During Development}
\label{sec:bugs}

We report three bugs discovered during implementation and
benchmarking, for transparency and because they represent failure
modes that other implementors may encounter.

\paragraph{Bug 1: entity\_ids field not populated in retrieval results.}
The initial Phase~2 implementation returned zero cross-stage results.
Root cause: the memory dicts returned by the PostgreSQL retrieval
query did not include the \texttt{entity\_ids} field (it was not in
the SELECT clause).  Fix: adding a LEFT JOIN to the
\texttt{entity\_memory} table and an \texttt{ARRAY\_AGG} of entity
IDs to the SELECT clause.

\paragraph{Bug 2: token budget constraining selection count.}
The initial submodular selection enforced the token budget as a hard
constraint: if a candidate's token count would exceed the remaining
budget, it was skipped.  This caused the assembler to select 2--3
items instead of 5.  Fix: decouple selection (by count) from budget
enforcement (at assembly time).

\paragraph{Bug 3: FlashRank silent failure.}
FlashRank's ONNX model occasionally fails to load and returns 0.0
scores for all candidates without raising an exception, producing
random rankings for both conditions.  Fix: a preflight check before
each benchmark run that reranks a synthetic pair and asserts non-zero
scores.

%----------------------------------------------------------------------
\section{Provenance and Prior Art}
\label{sec:provenance}

\subsection{Motivation for Provenance Documentation}

The structured context assembly architecture was designed before the
BEAM benchmark existed.  Because we report results on BEAM that
exceed previously published numbers (with the caveat that our metric
differs), it is important to document the chronological relationship
between the architecture's development and the benchmark's
publication.  We do so with verifiable commit SHAs from public
repositories.

\subsection{Timeline}

\begin{table}[t]
\centering
\caption{Development timeline with verifiable commit SHAs.}
\label{tab:timeline}
\small
\begin{tabular}{llll}
\toprule
Date & Event & Repo & SHA \\
\midrule
2025-09-14 & Context-aware PRD generation & ai-prd-builder & \texttt{3ef6c3f} \\
2025-09-25 & Modular arch.\ with cognitive modes & ai-prd-builder & \texttt{4f90564} \\
\textbf{2025-09-30} & \textbf{ContextManager.swift} & ai-prd-builder & \texttt{462de01} \\
2025-10-04 & Context-aware codebase interceptor & ai-prd-builder & \texttt{0743b0e} \\
\textbf{2025-10-31} & \textbf{BEAM paper published} & arXiv & --- \\
2025-07-10 & MIRIX published & arXiv & 2507.07957 \\
2026-02-10 & Hierarchical chunking + compression & ai-prd-generator & \texttt{8bc58f1} \\
2026-02-27 & ContextDecomposer & ai-architect* & \texttt{ba99681} \\
2026-03-03 & StageAwareContextAssembler & ai-architect* & \texttt{d4e2eb2} \\
2026-03-16 & Object-centric context decomposition & ai-architect* & \texttt{dc3c71d} \\
2026-04-07 & Python port + BEAM integration & Cortex & \texttt{5348f74} \\
\bottomrule
\multicolumn{4}{l}{\footnotesize *ai-architect-prd-builder (private repository)} \\
\end{tabular}
\end{table}

\subsection{The ContextManager.swift Origin}

The \texttt{ContextManager.swift} module, committed on September 30,
2025 in the public ai-prd-builder repository, implements the following
features that directly ancestor the ContextDecomposer:

\begin{itemize}
  \item Per-section token-budgeted context assembly with
    provider-specific limits (Apple Intelligence: 3,500 tokens,
    OpenAI: model-dependent)
  \item Slot-based budget splitting: 1/3 core request + 1/4
    clarifications + 1/6 tech stack, with the remaining budget
    distributed to lower-priority sections
  \item Section-keyword relevance filtering: only sections whose
    keywords match the current task are included
  \item Truncation awareness: the model is told which sections were
    reduced and by how much
\end{itemize}

The key design principle---``don't try to fit everything; structure
what goes in, prioritize it, and tell the model what was cut''---was
formulated for generating coherent 9-page PRDs on Apple Intelligence's
4,096-token context window.

\subsection{Evolution to Cortex}

The architecture evolved through three codebases:

\begin{enumerate}
  \item \textbf{ai-prd-builder} (September 2025): ContextManager with
    fixed-section budget splitting.  Public repository.

  \item \textbf{ai-architect-prd-builder} (February--March 2026):
    ContextDecomposer with priority-driven progressive condensation,
    domain-aware condensers, and truncation warning injection.
    StageAwareContextAssembler with two-phase 60/30 (with planned 10% summary tier) retrieval.
    Private repository.

  \item \textbf{Cortex} (April 2026): Python port adapted for
    Cortex's PostgreSQL-backed memory system, complemented with
    paper-backed mechanisms (HippoRAG PPR formulation, Krause \&
    Guestrin submodular guarantee).  Public repository.
\end{enumerate}

\subsection{Relationship to Published Systems}

The full combination has no published work we are aware of that combines these specific primitives in the 2024--2026
literature we surveyed across six research areas: computational
neuroscience, mathematics (submodularity, dimensionality reduction),
AI systems (RAG, long-context, agent memory), PhD theses on
computational memory models, vector database engineering, and
information theory.

The individual building blocks each have strong paper backing:

\begin{table}[h]
\centering
\caption{Paper backing for individual components.}
\label{tab:paper-backing}
\begin{tabular}{lll}
\toprule
Component & Paper & Year \\
\midrule
Submodular coverage & Krause \& Guestrin, JMLR & 2008 \\
MMR diversity & Carbonell \& Goldstein, SIGIR & 1998 \\
Personalized PageRank & Brin \& Page / Gutierrez et al. & 1998/2024 \\
Schema-structured summaries & Tse et al., Science & 2007 \\
Active retrieval & Wang \& Chen (MIRIX) & 2025 \\
RRF fusion & Cormack et al. & 2009 \\
\bottomrule
\end{tabular}
\end{table}

The contribution is the composition: the specific way these pieces
are combined into a two-primitive architecture that manages both
retrieval composition (StageAwareContextAssembler) and prompt-level
budgeting (ContextDecomposer).

%----------------------------------------------------------------------
\section{Limitations and Future Work}
\label{sec:limitations}

\subsection{Summarization Regression}

Summarization regresses at 10M scale ($-0.036$).  Stage-scoped
retrieval focuses depth at the cost of breadth.  Phase~3 (summary
fallback) is designed to address this but is currently a stub.
Wiring it to Cortex's CLS consolidation engine is the immediate next
step.  We predict elimination of the regression but have not verified
experimentally.

\subsection{Event Ordering Stagnation}

Event ordering scores 0.067 for both conditions---effectively random.
This ability requires precise chronological sequencing across
sessions, which is a temporal reasoning problem rather than a
retrieval problem.

\subsection{Entity Extraction Quality}

Entity extraction is rule-based: capitalized multi-word phrases,
quoted strings, technical terms with camelCase or snake\_case
patterns.  On BEAM's synthetic conversations (clear entity mentions,
consistent naming), this works adequately.  On real-world
conversational data with informal phrasing (``that Redis thing from
last week''), abbreviations (``the auth bug''), and implicit
references (``the same approach''), extraction quality would degrade
significantly.  Upgrading to a lightweight NER model (SpaCy,
fine-tuned token classifier) or LLM-based extraction is a natural
improvement.

\subsection{Latency Overhead}

\begin{table}[h]
\centering
\caption{Latency comparison (BEAM-10M).}
\label{tab:latency}
\begin{tabular}{lccc}
\toprule
Operation & WRRF & Assembler & Overhead \\
\midrule
Per-conversation (10M) & 307s & 733s & 2.4$\times$ \\
Per-query (10M) & $\sim$15s & $\sim$37s & 2.4$\times$ \\
\bottomrule
\end{tabular}
\end{table}

The overhead comes from per-query entity extraction from Phase~1
results, PPR graph construction and iteration, and submodular
selection's $O(nk)$ greedy loop.  Optimizations available but not
implemented include PPR precomputation, batched graph queries, and
lazy greedy acceleration \citep{Minoux1978}.

\subsection{End-to-End QA Evaluation}

All results are retrieval-proxy MRR.  We have not evaluated whether
improved retrieval translates to improved answer quality when a reader
model consumes the assembled context.  An LLM-as-judge evaluation
using BEAM's nugget scoring protocol would enable direct comparison
with LIGHT and other published systems.

\subsection{Budget Split Optimization}

The 60/30/10 split is a design heuristic inherited from the Swift PRD
generator, not an optimized parameter.  Systematic optimization (\eg,
Bayesian optimization over $(\beta_1, \beta_2, \beta_3)$ with
BEAM-10M MRR as the objective) could improve results.

\subsection{Stage Detection for Free-Form Conversations}

Our experiments use ground-truth plan IDs from BEAM.  Production
deployment requires robust stage detection without explicit labels.
The TemporalStageDetector (4-hour gap threshold) is reasonable but
untested on BEAM.

\subsection{Larger Embedding Models}

Our experiments use all-MiniLM-L6-v2 (384D).  The JL lower bound
analysis suggests that 768D or 1024D embeddings could push the
geometric ceiling higher, reducing the need for structured assembly.
An important open question: does the assembly architecture remain
beneficial with state-of-the-art 1024D models (\eg, E5-Mistral,
NV-Embed-2), or does the improved embedding space make flat retrieval
sufficient even at 10M tokens?

\subsection{Generalization Beyond Conversational Memory}

The architecture was designed for and tested on conversational memory.
Whether it generalizes to other long-context retrieval
problems---legal document search, medical record retrieval, codebase
navigation---is unknown.  The stage-detection mechanism is generic
(any partitioning of the corpus), but the specific condensers and
entity extraction are tuned for conversation.

%----------------------------------------------------------------------
\section{Conclusion}
\label{sec:conclusion}

We presented a structured context assembly architecture for long-term
memory retrieval that addresses the fundamental scaling limitation of
dense vector retrieval: at sufficient corpus density, embedding-based
nearest-neighbor search cannot discriminate relevant from irrelevant
memories because the geometric structure of the embedding space
degenerates.

The architecture introduces two composable primitives.  The
\textbf{ContextDecomposer} manages prompt-level budgeting with typed
priority slots, domain-aware condensers, and an explicit truncation
warning mechanism that communicates assembly decisions to the reader
model---a mechanism with no published work we are aware of that combines these specific primitives in the retrieval or
prompt engineering literature.  The \textbf{StageAwareContextAssembler}
manages retrieval-level composition through three phases: own-stage
retrieval with submodular diversity \citep{Krause2008}, cross-stage
bridging via Personalized PageRank over the entity graph
\citep{Gutierrez2024}, and summary fallback for uncovered stages
\citep{Tse2007}.

The empirical results on BEAM validate the scale-dependent thesis:
the architecture is net-flat at 100K tokens ($+1.9\%$, within noise)
and provides a $+21.5\%$ improvement at 10M tokens, with 8 of 10
memory abilities improving.  The multi-session reasoning sign flip
(from $-0.312$ at 100K to $+0.128$ at 10M) is the strongest evidence
that the benefit is genuinely scale-dependent rather than incidental.

The architecture does not introduce novel retrieval primitives.
Every building block has strong paper backing.  The contribution is
the composition: priority budgeting, stage-aware three-phase assembly,
domain-aware condensation, and model-facing truncation awareness,
combined into a system that manages context at a level of abstraction
above individual retrieval signals.

The insight that motivated this work---that at scale, \emph{what}
enters context matters more than \emph{how well} individual items are
retrieved---emerged from a practical constraint: generating coherent
documents on a 4,096-token context window.  That the same
architectural principle applies to 10-million-token memory retrieval
suggests it may be general: whenever the available context is smaller
than the relevant information, structured assembly outperforms flat
ranking.

The crossover point between flat retrieval and structured assembly is
empirically located between 100K and 10M tokens for 384D embeddings.
As embedding models improve and context windows grow, this crossover
point will shift---but the fundamental principle will remain: there
will always be a scale at which the relevant information exceeds the
available context, and at that scale, structured assembly will
dominate.

%----------------------------------------------------------------------
\section*{Acknowledgments}

The Python port, benchmark integration, and paper-backed complements
were implemented with Claude Opus 4.6 (Anthropic).

%----------------------------------------------------------------------
\appendix

\section{Reproducing the Experiments}
\label{app:reproducing}

\subsection{Prerequisites}

PostgreSQL 17+ with pgvector and pg\_trgm extensions, Python 3.10+,
and Cortex installed from source:

\begin{verbatim}
git clone https://github.com/cdeust/Cortex.git
cd Cortex
pip install -e ".[postgresql,benchmarks]"
\end{verbatim}

\subsection{Running Benchmarks}

\begin{verbatim}
# BEAM-100K baseline (WRRF only)
DATABASE_URL="postgresql://localhost:5432/cortex_bench" \
  python benchmarks/beam/run_benchmark.py --split 100K

# BEAM-10M baseline (WRRF only)
DATABASE_URL="postgresql://localhost:5432/cortex_bench" \
  python benchmarks/beam/run_benchmark.py --split 10M

# BEAM-10M with structured context assembly
CORTEX_USE_ASSEMBLER=1 \
DATABASE_URL="postgresql://localhost:5432/cortex_bench" \
  python benchmarks/beam/run_benchmark.py --split 10M

# BEAM-100K with structured context assembly
CORTEX_USE_ASSEMBLER=1 \
DATABASE_URL="postgresql://localhost:5432/cortex_bench" \
  python benchmarks/beam/run_benchmark.py --split 100K
\end{verbatim}

\subsection{Expected Runtimes}

\begin{table}[h]
\centering
\caption{Expected benchmark runtimes.}
\label{tab:runtimes}
\begin{tabular}{lccc}
\toprule
Configuration & Conversations & Time/Conv & Total \\
\midrule
BEAM-100K WRRF & 5 & $\sim$32s & $\sim$2.5 min \\
BEAM-100K Assembler & 5 & $\sim$293s & $\sim$24 min \\
BEAM-10M WRRF & 10 & $\sim$307s & $\sim$51 min \\
BEAM-10M Assembler & 10 & $\sim$733s & $\sim$122 min \\
\bottomrule
\end{tabular}
\end{table}

\subsection{Database Protocol}

The benchmark harness creates a fresh \texttt{cortex\_bench} database
per run.  Between conversations within a run, all data tables are
truncated:

\begin{verbatim}
TRUNCATE memories, entities, relationships,
  entity_memory CASCADE;
\end{verbatim}

This ensures zero cross-conversation contamination.  The TRUNCATE
approach (vs.\ DROP/CREATE) preserves the schema, extensions, and
stored procedures, reducing per-conversation setup time.

\subsection{Variance Script}

The full variance analysis is run via:

\begin{verbatim}
bash benchmarks/beam/variance/bench_variance.sh
\end{verbatim}

This script runs each configuration 3 times, with PostgreSQL restart
between runs, and reports mean and standard deviation of MRR.
Observed variance: $\pm 0.01$ MRR.

%----------------------------------------------------------------------
\section{Full Ablation Results}
\label{app:ablation}

\subsection{Per-Ability Ablation on BEAM-100K (5 Conversations)}

See Table~\ref{tab:ablation-detail} in Section~\ref{sec:ablation} for
the full per-ability breakdown.

\subsection{Configuration Descriptions}

\begin{itemize}
  \item \textbf{WRRF}: flat retrieval over the entire corpus, no
    assembler.  Standard top-$k$ selection after WRRF fusion +
    FlashRank reranking.

  \item \textbf{No Submod}: assembler enabled with stage-scoping, but
    Phase~1 uses naive top-$k$ selection instead of submodular
    coverage.  Isolates the effect of stage-scoping without diversity.

  \item \textbf{Ph1 Only}: assembler with submodular coverage
    selection in Phase~1, but Phase~2 (PPR traversal) disabled.

  \item \textbf{Ph1+Ph2}: assembler with submodular coverage in
    Phase~1 and PPR traversal in Phase~2.  Phase~3 is a stub.

  \item \textbf{Full}: same as Ph1+Ph2 in the current implementation
    (Phase~3 not yet wired to real summaries).
\end{itemize}

\subsection{Key Ablation Findings}

\begin{enumerate}
  \item \textbf{Stage-scoping without diversity hurts} ($-0.079$ vs.\
    WRRF).  The No~Submod configuration reduces the candidate pool to
    one stage's memories and then picks the top-$k$ by score, which
    selects redundant memories.

  \item \textbf{Submodular selection alone (Ph1 Only) is not
    sufficient} ($-0.078$ vs.\ WRRF).  Without Phase~2 PPR bridging,
    the assembler is limited to same-stage memories.

  \item \textbf{The combination of submodular + PPR recovers and
    slightly exceeds WRRF} ($+0.011$).  Neither mechanism alone helps
    at 100K; together, they produce a marginal improvement.

  \item \textbf{At 10M, we expect both effects to be amplified}, based
    on the scale-dependent analysis in Section~\ref{sec:scale-dependent}.
\end{enumerate}

%----------------------------------------------------------------------
\section{Implementation Architecture}
\label{app:implementation}

\subsection{Module Structure}

The context assembly module resides in
\texttt{mcp\_server/core/context\_assembly/} within the Cortex
codebase:

\begin{table}[h]
\centering
\caption{Module structure of the context assembly implementation.}
\label{tab:modules}
\small
\begin{tabular}{lrcl}
\toprule
Module & Lines & Purpose & Paper Backing \\
\midrule
\texttt{decomposer.py} & 195 & Priority-budgeted assembly & Novel \\
\texttt{stage\_assembler.py} & 320 & Three-phase assembler & Novel composition \\
\texttt{coverage.py} & 128 & Submodular selection (Ph1) & Krause \& Guestrin \\
\texttt{ppr\_traversal.py} & 177 & PPR traversal (Ph2) & Gutierrez et al. \\
\texttt{stage\_detector.py} & 196 & Pluggable stage detection & Novel interface \\
\texttt{condensers.py} & 277 & Domain-aware condensers & Novel \\
\texttt{budget.py} & 123 & Token estimation/allocation & Ported from Swift \\
\texttt{warning.py} & 62 & Truncation warning banner & Novel \\
\texttt{active\_retrieval.py} & 189 & Query reformulation & Wang \& Chen \\
\bottomrule
\multicolumn{4}{l}{\footnotesize Total: $\sim$1,667 lines of Python across 10 modules.} \\
\end{tabular}
\end{table}

\subsection{Architectural Invariants}

All modules are in the \texttt{core/} layer (pure business logic, zero
I/O).  This is enforced by Cortex's clean architecture:
\texttt{core/} modules import only from \texttt{shared/} and Python
stdlib.  External dependencies (PostgreSQL queries, embedding
generation, schema engine) are injected as callbacks at construction
time.  Handlers (the composition root layer) wire \texttt{core/} to
\texttt{infrastructure/} via dependency injection.

This means the context assembly logic can be tested entirely in memory
with mock callbacks, without a running database.  The test suite
includes 47 tests covering all modules.

\subsection{Callback Interface}

The \texttt{StageAwareContextAssembler} receives five callbacks:

\begin{itemize}
  \item \texttt{retrieve\_fn(query, stage\_id, max\_results)}: Phase~1
    retrieval.  Wired to the WRRF pipeline with a stage filter.
  \item \texttt{entity\_graph\_fn()}: returns (entities,
    relationships) for PPR.  Wired to PostgreSQL entity/relationship
    queries.
  \item \texttt{memories\_by\_entity\_fn(entity\_ids)}: returns
    memories containing the given entities.  Wired to the
    \texttt{entity\_memory} junction table.
  \item \texttt{stage\_summary\_fn(stage\_id)}: returns a summary for
    the given stage.  Currently a stub returning empty string.
\end{itemize}

This callback design ensures that the assembler is agnostic to the
underlying storage and retrieval infrastructure.

%----------------------------------------------------------------------
\section{BEAM-100K Full Results (20 Conversations, WRRF Control)}
\label{app:100k-full}

For completeness, we include the full WRRF results on the BEAM-100K
20-conversation split (all conversations, not just the 5 used for
A/B comparison):

\begin{table}[h]
\centering
\caption{BEAM-100K full WRRF results (20 conversations).}
\label{tab:100k-full}
\begin{tabular}{lcccc}
\toprule
Ability & MRR & R@5 & R@10 & Qs \\
\midrule
contradiction\_resolution & 0.729 & 90.0\% & 90.0\% & 40 \\
temporal\_reasoning & 0.684 & 75.0\% & 80.0\% & 40 \\
knowledge\_update & 0.735 & 87.5\% & 90.0\% & 40 \\
multi\_session\_reasoning & 0.596 & 77.5\% & 80.0\% & 40 \\
event\_ordering & 0.311 & 50.0\% & 62.5\% & 40 \\
summarization & 0.315 & 44.4\% & 61.1\% & 36 \\
preference\_following & 0.309 & 52.5\% & 60.0\% & 39 \\
instruction\_following & 0.219 & 27.5\% & 42.5\% & 40 \\
abstention & 0.125 & 12.5\% & 12.5\% & 40 \\
\midrule
\textbf{Overall} & \textbf{0.438} & --- & --- & 395 \\
\bottomrule
\end{tabular}
\end{table}

This 20-conversation result (0.438 MRR) is lower than the
5-conversation result (0.591 MRR), reflecting variance across
conversations.  The 5 conversations selected for A/B comparison
happen to be above-average in difficulty.

%----------------------------------------------------------------------
\section{Development Timeline and Intermediate Results}
\label{app:dev-timeline}

\subsection{Intermediate BEAM-10M Runs}

During development, several intermediate assembler configurations
were benchmarked on BEAM-10M:

\begin{table}[h]
\centering
\caption{Intermediate BEAM-10M results during development.}
\label{tab:intermediate}
\begin{tabular}{lcp{5cm}}
\toprule
Configuration & MRR & Issue \\
\midrule
First run (entity\_ids bug) & 0.100 & Phase~2 returned 0 results \\
After entity\_ids fix & 0.120 & Stage detection used local turn IDs \\
After stage fix (final) & 0.429 & Correct stage assignment \\
\bottomrule
\end{tabular}
\end{table}

The jump from 0.120 to 0.429 upon fixing the stage detection confirms
that correct stage assignment is critical to the architecture's
performance.

\subsection{Debugging the entity\_ids Bug}

The initial Phase~2 implementation consistently returned zero
cross-stage memories.  The debugging trace: (1)~PPR returned non-zero
scores for entities, but \texttt{memories\_by\_entity\_fn} returned
empty lists.  (2)~The retrieval query's SELECT clause did not include
\texttt{entity\_ids}.  (3)~The field was populated at ingest time but
not JOINed at retrieval time.  (4)~Fix: add LEFT JOIN and
ARRAY\_AGG.

\subsection{Debugging the Stage Detection Bug}

BEAM-10M turn IDs are plan-relative (each plan restarts numbering
from 0), but gold annotations reference global turn indices.  Without
cumulative plan offsets, the stage detector assigned all memories to
``plan-0.''  Fix: compute cumulative turn offsets per plan and apply
before gold matching.

%----------------------------------------------------------------------
\section{Pseudocode for Core Algorithms}
\label{app:pseudocode}

\subsection{assemble\_prompt (ContextDecomposer)}

\begin{algorithm}[H]
\caption{ContextDecomposer prompt assembly}
\label{alg:assemble-prompt}
\begin{algorithmic}[1]
\REQUIRE template, placeholders, context\_window, headroom
\STATE budget $\leftarrow \lfloor$ context\_window $\times$ headroom $\rfloor$
\STATE shell $\leftarrow$ substitute(template, $\{p.\text{key}: \text{``''} \mid p \in \text{placeholders}\}$)
\STATE variable\_budget $\leftarrow \max(300, \text{budget} - \tau(\text{shell}))$

\STATE \COMMENT{Fast path}
\IF{$\sum_p \tau(p.\text{value}) \leq$ variable\_budget}
  \RETURN fill\_template(template, placeholders)
\ENDIF

\STATE \COMMENT{Progressive condensation}
\STATE sorted\_ph $\leftarrow$ sort(placeholders, by priority DESC)
\STATE remaining $\leftarrow$ variable\_budget
\STATE effective $\leftarrow \{\}$
\FOR{$i, p$ \textbf{in} enumerate(sorted\_ph)}
  \STATE not\_yet $\leftarrow |\text{sorted\_ph}| - i$
  \STATE share $\leftarrow \max(50, \lfloor \text{remaining} / \max(1, \text{not\_yet}) \rfloor)$
  \IF{$\tau(p.\text{value}) \leq$ share}
    \STATE effective[$p.\text{key}$] $\leftarrow p.\text{value}$
    \STATE remaining $\leftarrow$ remaining $- \tau(p.\text{value})$
  \ELSE
    \IF{$p.\text{condenser} \neq$ null}
      \STATE reduced $\leftarrow p.\text{condenser}(p.\text{value}, \text{share})$
    \ELSE
      \STATE reduced $\leftarrow$ truncate(${p.\text{value}}$, share)
    \ENDIF
    \STATE effective[$p.\text{key}$] $\leftarrow$ reduced
    \STATE remaining $\leftarrow$ remaining $- \min(\tau(\text{reduced}), \text{remaining})$
  \ENDIF
\ENDFOR

\STATE \COMMENT{Post-assembly safety}
\STATE prompt $\leftarrow$ fill\_template(template, effective)
\WHILE{$\tau(\text{prompt}) >$ context\_window $-$ safety\_margin}
  \STATE halve\_lowest\_priority(effective, sorted\_ph)
  \STATE prompt $\leftarrow$ fill\_template(template, effective)
\ENDWHILE

\STATE \COMMENT{Truncation warning}
\STATE banner $\leftarrow$ build\_truncation\_banner(original, final)
\IF{banner $\neq$ ``''}
  \STATE prompt $\leftarrow$ banner $+$ ``\textbackslash n\textbackslash n'' $+$ prompt
\ENDIF
\RETURN prompt
\end{algorithmic}
\end{algorithm}

\subsection{assemble (StageAwareContextAssembler)}

\begin{algorithm}[H]
\caption{StageAwareContextAssembler}
\label{alg:assemble}
\begin{algorithmic}[1]
\REQUIRE query, current\_stage, budget\_split, max\_chunks
\STATE \COMMENT{Phase 1: Own-stage}
\STATE candidates $\leftarrow$ retrieve(query, current\_stage, max\_chunks $\times$ 3)
\STATE phase1 $\leftarrow$ submodular\_select(candidates, max\_chunks, $\lambda = 0.5$)

\STATE \COMMENT{Phase 2: Cross-stage via PPR}
\STATE seed\_entities $\leftarrow$ extract\_entities(phase1)
\STATE entities, relationships $\leftarrow$ entity\_graph()
\STATE adjacency $\leftarrow$ build\_adjacency(entities, relationships)
\STATE ppr $\leftarrow$ PPR(adjacency, seed\_entities, $\alpha = 0.15$)
\STATE top\_entities $\leftarrow$ top\_$k$(ppr, 50)
\STATE cross $\leftarrow$ memories\_by\_entities(top\_entities)
\STATE cross $\leftarrow$ filter(cross, stage $\neq$ current\_stage)
\STATE phase2 $\leftarrow$ score\_by\_ppr(cross, ppr)[:max\_chunks]

\STATE \COMMENT{Phase 3: Summary fallback}
\STATE covered $\leftarrow \{$current\_stage$\} \cup$ stages\_of(phase2)
\STATE uncovered $\leftarrow$ all\_stages $-$ covered
\STATE phase3 $\leftarrow$ [stage\_summary($s$) for $s$ in uncovered]

\RETURN structured\_context(phase1, phase2, phase3)
\end{algorithmic}
\end{algorithm}

%----------------------------------------------------------------------

\bibliographystyle{plainnat}
\bibliography{references}

\end{document}
